\section{AIME II 2017}
        \begin{enumerate}
        	\item (\textit{AIME II 2017}) Find the number of subsets of $\{1, 2, 3, 4, 5, 6, 7, 8\}$ that are subsets of neither $\{1, 2, 3, 4, 5\}$ nor $\{4, 5, 6, 7, 8\}$.


 
 \textbf{Answer: }196



\textbf{Solution 1}

The number of subsets of a set with $n$ elements is $2^n$. The total number of subsets of $\{1, 2, 3, 4, 5, 6, 7, 8\}$ is equal to $2^8$. The number of sets that are subsets of at least one of $\{1, 2, 3, 4, 5\}$ or $\{4, 5, 6, 7, 8\}$ can be found using complementary counting. There are $2^5$ subsets of $\{1, 2, 3, 4, 5\}$ and $2^5$ subsets of $\{4, 5, 6, 7, 8\}$. It is easy to make the mistake of assuming there are $2^5+2^5$ sets that are subsets of at least one of $\{1, 2, 3, 4, 5\}$ or $\{4, 5, 6, 7, 8\}$, but the $2^2$ subsets of $\{4, 5\}$ are overcounted. There are $2^5+2^5-2^2$ sets that are subsets of at least one of $\{1, 2, 3, 4, 5\}$ or $\{4, 5, 6, 7, 8\}$, so there are $2^8-(2^5+2^5-2^2)$ subsets of $\{1, 2, 3, 4, 5, 6, 7, 8\}$ that are subsets of neither $\{1, 2, 3, 4, 5\}$ nor $\{4, 5, 6, 7, 8\}$. $2^8-(2^5+2^5-2^2)=\boxed{196}$.

 





\textbf{Solution 2}

Upon inspection, a viable set must contain at least one element from both of the sets $\{1, 2, 3, 4, 5\}$ and $\{4, 5, 6, 7, 8\}$. Since 4 and 5 are included in both of these sets, then they basically don't matter, i.e. if set A is a subset of both of those two then adding a 4 or a 5 won't change that fact. Thus, we can count the number of ways to choose at least one number from 1 to 3 and at least one number from 6 to 8, and then multiply that by the number of ways to add in 4 and 5. The number of subsets of a 3 element set is $2^3=8$, but we want to exclude the empty set, giving us 7 ways to choose from  $\{1, 2, 3\}$ or $\{6, 7, 8\}$. We can take each of these $7 \times 7=49$ sets and add in a 4 and/or a 5, which can be done in 4 different ways (by adding both, none, one, or the other one). Thus, the answer is $49 \times 4=\boxed{196}$.

 





\textbf{Solution 3}

This solution is very similar to Solution $2$. The set of all subsets of $\{1,2,3,4,5,6,7,8\}$ that are disjoint with respect to $\{4,5\}$ and are not disjoint with respect to the complements of sets (and therefore not a subset of) $\{1,2,3,4,5\}$ and $\{4,5,6,7,8\}$ will be named $S$, which has $7\cdot7=49$ members. The union of each member in $S$ and the $2^2=4$ subsets of $\{4,5\}$ will be the members of set $Z$, which has $49\cdot4=\boxed{196}$ members. $\blacksquare$

 Solution by \href{https://artofproblemsolving.com/wiki/index.php?title=User:A1b2}{a1b2}

 





\textbf{Solution 4}

Consider that we are trying to figure out how many subsets are possible of $\{1,2,3,4,5,6,7,8\}$ that are not in violation of the two subsets $\{1,2,3,4,5\}$ and $\{4,5,6,7,8\}$. Assume that the number of numbers we pick from the subset $\{1,2,3,4,5,6,7,8\}$ is $n$. Thus, we can compute this problem with simple combinatorics:

 If $n=1$, $\binom{8}{1}$ $-$ ($\binom{5}{1}$ + $\binom{5}{1}$ $-  2$) [subtract $2$ to eliminate the overcounting of the subset $\{4\}$ or $\{5\}$] = $8 - 8$ = $0$

 If $n=2$, $\binom{8}{2}$ $-$ ($\binom{5}{2}$ + $\binom{5}{2}$ $-  1$) [subtract $1$ to eliminate the overcounting of the subset $\{4,5\}$] = $28 - 19$ = $9$

 If $n=3$, $\binom{8}{3}$ $-$ ($\binom{5}{3}$ + $\binom{5}{3}$) = $56 - 20$ = $36$

 If $n=4$, $\binom{8}{4}$ $-$ ($\binom{5}{4}$ + $\binom{5}{4}$) = $70 - 10$ = $60$

 If $n=5$, $\binom{8}{5}$ $-$ ($\binom{5}{5}$ + $\binom{5}{5}$) = $56 - 2$ = $54$

 If $n>5$, then the set $\{1,2,3,4,5,6,7,8\}$ is never in violation of the two subsets $\{1,2,3,4,5\}$ and $\{4,5,6,7,8\}$. Thus,

 If $n=6$, $\binom{8}{6}$ = $28$

 If $n=7$, $\binom{8}{7}$ = $8$

 If $n=8$, $\binom{8}{8}$ = $1$

 Adding these together, our solution becomes $0$ + $9$ + $36$ + $60$ + $54$ + $28$ + $8$ + $1$ $=\boxed{196}$

 Solution by IronicNinja~

 





\textbf{Solution 5 (Solution 1, but a little different, using PIE)}

PIE is the principle of inclusion and exclusion. Denote $\{1, 2, 3, 4, 5, 6, 7, 8\}$ as set $T$, $\{1, 2, 3, 4, 5\}$ as set $A$, and $\{4, 5, 6, 7, 8\}$ as set $B$. We want to find the number of subsets of $T$ that are not subsets of $A$ or $B$ (i.e., subsets contained in $T$ but not in $A \cup B$). Hence, the desired number is $|T| - |A \cup B|$.

 By PIE, $|A \cup B| = |A| + |B| - |A \cap B|$. Notice that $A \cap B = \{4, 5\}$, so $|A \cap B| = 4$ (the number of subsets of $\{4,5\}$). Both $A$ and $B$ contain $2^5 = 32$ subsets, and $T$ contains $2^8 = 256$ subsets. The final answer is therefore
 \[256 - (32 + 32 - 4) = 256 - 60 = \boxed{196}.\]

 




	\item (\textit{AIME II 2017}) Teams $T_1$, $T_2$, $T_3$, and $T_4$ are in the playoffs. In the semifinal matches, $T_1$ plays $T_4$, and $T_2$ plays $T_3$. The winners of those two matches will play each other in the final match to determine the champion. When $T_i$ plays $T_j$, the probability that $T_i$ wins is $\frac{i}{i+j}$, and the outcomes of all the matches are independent. The probability that $T_4$ will be the champion is $\frac{p}{q}$, where $p$ and $q$ are relatively prime positive integers. Find $p+q$.


 
 \textbf{Answer: }781



\textbf{Solution 1}

There are two scenarios in which $T_4$ wins. The first scenario is where $T_4$ beats $T_1$, $T_3$ beats $T_2$, and $T_4$ beats $T_3$, and the second scenario is where $T_4$ beats $T_1$, $T_2$ beats $T_3$, and $T_4$ beats $T_2$. Consider the first scenario. The probability $T_4$ beats $T_1$ is $\frac{4}{4+1}$, the probability $T_3$ beats $T_2$ is $\frac{3}{3+2}$, and the probability $T_4$ beats $T_3$ is $\frac{4}{4+3}$. Therefore the first scenario happens with probability $\frac{4}{4+1}\cdot\frac{3}{3+2}\cdot\frac{4}{4+3}$. Consider the second scenario. The probability $T_4$ beats $T_1$ is $\frac{4}{1+4}$, the probability $T_2$ beats $T_3$ is $\frac{2}{2+3}$, and the probability $T_4$ beats $T_2$ is $\frac{4}{4+2}$. Therefore the second scenario happens with probability $\frac{4}{1+4}\cdot\frac{2}{2+3}\cdot\frac{4}{4+2}$. By summing these two probabilities, the probability that $T_4$ wins is $\frac{4}{4+1}\cdot\frac{3}{3+2}\cdot\frac{4}{4+3}+\frac{4}{1+4}\cdot\frac{2}{2+3}\cdot\frac{4}{4+2}$. Because this expression is equal to $\frac{256}{525}$, the answer is $256+525=\boxed{781}$.

 





\textbf{Solution 2}

Clearly $T_4$ has to win its game with $T_1$, which has probability $\frac{4}{5}$. There are two cases, depending on who its opponent is.Case 1: $T_4$ faces $T_2$. So $T_2$ won its first game with probability $\frac{2}{5}$, and $T_4$ wins the finals with probability $\frac{4}{6}=\frac{2}{3}$.Case 2: $T_4$ faces $T_3$. So $T_3$ won its first game with probability $\frac{3}{5}$, and $T_4$ wins the finals with probability $\frac{4}{7}$.

 The total probability is therefore $\frac{4}{5} \times \left(\frac{4}{15} + \frac{12}{35}\right) = \frac{4}{5} \cdot \frac{64}{105} = \frac{256}{525} \implies 256+525 = \boxed{781}$.~First

 




	\item (\textit{AIME II 2017}) A triangle has vertices $A(0,0)$, $B(12,0)$, and $C(8,10)$. The probability that a randomly chosen point inside the triangle is closer to vertex $B$ than to either vertex $A$ or vertex $C$ can be written as $\frac{p}{q}$, where $p$ and $q$ are relatively prime positive integers. Find $p+q$.


 
 \textbf{Answer: }409



\textbf{Solution 1}


                    \begin{figure}[H]
                        \centering
                        \includegraphics[width=0.5\textwidth]{Solution/AIME_2017_II_3_Solution_Figure_1.png}
                    \end{figure}
                    

 The set of all points closer to point $B$ than to point $A$ lie to the right of the perpendicular bisector of $AB$ (line $PZ$ in the diagram), and the set of all points closer to point $B$ than to point $C$ lie below the perpendicular bisector of $BC$ (line $PX$ in the diagram). Therefore, the set of points inside the triangle that are closer to $B$ than to either vertex $A$ or vertex $C$ is bounded by quadrilateral $BXPZ$. Because $X$ is the midpoint of $BC$ and $Z$ is the midpoint of $AB$, $X=(10,5)$ and $Z=(6,0)$. The coordinates of point $P$ is the solution to the system of equations defined by lines $PX$ and $PZ$. Using the point-slope form of a linear equation and the fact that the slope of the line perpendicular to a line with slope $m$ is $-\frac{1}{m}$, the equation for line $PX$ is $y=\frac{2}{5}x+1$ and the equation for line $PZ$ is $x=6$. The solution of this system is $P=\left(6,\frac{17}{5}\right)$. Using the shoelace formula on quadrilateral $BXPZ$ and triangle $ABC$, the area of quadrilateral $BXPZ$ is $\frac{109}{5}$ and the area of triangle $ABC$ is $60$. Finally, the probability that a randomly chosen point inside the triangle is closer to vertex $B$ than to vertex $A$ or vertex $C$ is the ratio of the area of quadrilateral $BXPZ$ to the area of $ABC$, which is $\frac{\frac{109}{5}}{60}=\frac{109}{300}$. The answer is $109+300=\boxed{409}$.

 





\textbf{Solution 2}

Since we know the coordinates of all three vertices of the triangle, we can find the side lengths: $AB=12$, $AC=2\sqrt{41}$, and $BC=2\sqrt{29}$. We notice that the point where the three distances are the same is the circumcenter - so we use one of the triangle area formulas to find the circumradius, since we know what the area is.
 \[\frac{12 \cdot 2\sqrt{41} \cdot 2\sqrt{29}}{4 \cdot R}=\frac{12 \cdot 10}{2}.\]We rearrange to get
 \[R=\frac{\sqrt{41} \cdot \sqrt{29}}{5}.\]
                    \begin{figure}[H]
                        \centering
                        \includegraphics[width=0.5\textwidth]{Solution/AIME_2017_II_3_Solution_Figure_2.png}
                    \end{figure}
                    We know that $AP=\frac{\sqrt{41} \cdot \sqrt{29}}{5}$, and $AS=6$, so using the Pythagorean Theorem gives $SP=\frac{17}{5}$. This means $[ASP]=[BSP]=\frac{17}{5} \cdot 6 \cdot \frac{1}{2} = \frac{51}{5}$.Similarly, we know that $BP=\frac{\sqrt{41} \cdot \sqrt{29}}{5}$, and $BR=\sqrt{29}$, so we get that $PR=\frac{4\sqrt{29}}{5}$, and so $[BRP]=[CRP]=\frac{4\sqrt{29}}{5} \cdot \sqrt{29} \cdot \frac{1}{2} = \frac{58}{5}$.Lastly, we know that $CP=\frac{\sqrt{41} \cdot \sqrt{29}}{5}$, and $CT=\sqrt{41}$, so we get that $PT=\frac{2\sqrt{41}}{5}$, and $[ATP]=[CTP]=\frac{2\sqrt{41}}{5} \cdot \sqrt{41} \cdot \frac{1}{2} = \frac{41}{5}$.Therefore, our answer is $\frac{51+58}{2(51+58+41)}=\frac{109}{300} \rightarrow \boxed{409}$.

 





\textbf{Solution 3}

To start the problem, identify the two midpoints that connect $AB$ and $BC$. This is because the midpoints of such lines is the mark at which the point will sway closer to vertex $A$/$C$ or vertex $B$. The midpoint of $AB$ is $(6,0)$, and the midpoint of $BC$ is $(10,5)$. Then,  determine the line at which the distance between vertex $B$ and vertex $C$ are the same. Assuming that $x$ is the real value of $x$ and $y$ is the real value of $y$, we can create a simple equation:

 $\sqrt{|x-8|^2 + (10-y)^2}$ = $\sqrt{(12-x)^2 + y^2}$

 where the left side of the equation is for the distance to vertex $C$ and the right side of the equation is the distance to vertex $B$.

 Squaring both sides and then distributing, we get 

 $x^2 - 16x + 64 + y^2 - 20y + 100 = x^2 - 24x + 144 + y^2$.

 Notice that $(x-8)^2 = (-x+8)^2$, and thus there is no need to create another equation.

 Simplifying, we get $8x + 20 = 20y$.

 Divide both sides by 20, then simplify, and the line that represents equivalent distance between vertex $B$ and vertex $C$ is $y = \frac{2x}{5} + 1$.

 This line starts at the midpoint of $BC$, which is $(10,5)$, and ends at the line $x=6$, as $x=6$ represents equivalent distance between vertex $A$ and vertex $B$. Plug in $x=6$ to the equation $y = \frac{2x}{5} + 1$, and we get $y = \frac{17}{5}$. Now that we have our four points that are $(6,0), (6,\frac{17}{5}), (10,5)$, and $(12,0)$, we can calculate the area of the quadrilateral in which a point is closer to vertex $B$ as opposed to either vertex $A$ or vertex $C$. Simply draw a rectangle that has the points $(6,0), (6,5), (12,5)$ and $(12,0)$, and then subtract the two triangles that appear in between. 

 Thus, the area of the quadrilateral is $6\cdot5 - (\frac{\frac{8}{5}\cdot4}{2} + \frac{5\cdot2}{2}) \rightarrow 30 - \frac{41}{5} = \frac{109}{5}$. Since the problem asks us for the probability that a point chosen inside the triangle is inside the quadrilateral, and because the area of $\triangle ABC$ is $\frac{12\cdot10}{2} = 60$, the probability is $\frac{\frac{109}{5}}{60}=\frac{109}{300}$, which means the final answer is $109+300=\boxed{409}$.

 Solution by IronicNinja~

 





\textbf{Solution 4 (Shoelace Theorem/Formula)}

Calculate the area of the triangle using the \href{https://artofproblemsolving.com/wiki/index.php?title=Shoelace\_Theorem}{Shoelace Theorem} on $(0,0), (12,0), (8,10)$
 \[\frac{1}{2}|(0+120+0)-(0+0+0)|=60\]

 Get the four points $(6,0), (6,\frac{17}{5}), (10,5)$, and $(12,0)$ by any method from the above solutions. Then use the \href{https://artofproblemsolving.com/wiki/index.php?title=Shoelace\_Theorem}{Shoelace Theorem} to find the area of the region we want:
 \[\frac{1}{2}|(0+60+34+0)-(0+0+30+\frac{102}{5})|=\frac{109}{5}\]

 Therefore the probability is $\frac{\frac{109}{5}}{60}=\frac{109}{300}$. Thus giving the final answer of $109+300=\boxed{409}$.

 Solution by phoenixfire

 





\textbf{Solution 5}

Draw the circumradii from the circumcenter to the three vertices. Drop perpendicular from the circumcenter to the sides. Note that since the triangle is isosceles, the perpendicular are in fact perpendicular bisectors. Therefore the region containing the points closer to B are in $\triangle{OBP_{1}} , \triangle{OBP_{2}}$ where $O$ is the circumcenter and $P_{1}, P_{2}$ points of contact of the perpendiculars and the sides. Therefore our fraction is $\frac{109}{300}$. Our answer is then $\boxed{409}$.

 





\textbf{Solution 6}

Using the same graph and methods as Solution 1, find the coordinates of $P$, $X$, $B$, and $Z$. Also, note that angles $PXB$ and $PZB$ are right angles, so $PXBZ$ is a cyclic quadrilateral. Then, use Brahmagupta's formula to determine the area of the quadrilateral, which is $\frac{109}{5}$. Then find the area of triangle $ABC$, which is $60$. The probability is $\frac{109}{300}$. Finally, we get our answer of $\boxed{409}$

 




	\item (\textit{AIME II 2017}) Find the number of positive integers less than or equal to $2017$ whose base-three representation contains no digit equal to $0$.


 
 \textbf{Answer: }222



\textbf{Solution 1}

The base-$3$ representation of $2017_{10}$ is $2202201_3$. Because any $7$-digit base-$3$ number that starts with $22$ and has no digit equal to $0$ must be greater than $2017_{10}$, all $7$-digit numbers that have no digit equal to $0$ must start with $21$ or $1$ in base $3$. Of the base-$3$ numbers that have no digit equal to $0$, there are $2^5$ $7$-digit numbers that start with $21$, $2^6$ $7$-digit numbers that start with $1$, $2^6$ $6$-digit numbers, $2^5$ $5$-digit numbers, $2^4$ $4$-digit numbers, $2^3$ $3$-digit numbers, $2^2$ $2$-digit numbers, and $2^1$ $1$-digit numbers. Summing these up, we find that the answer is $2^5+2^6+2^6+2^5+2^4+2^3+2^2+2^1=\boxed{222}$.

 





\textbf{Solution 2}

Note that $2017=2202201_{3}$, and $2187=3^7=10000000_{3}$. There can be a $1,2,...,7$ digit number less than $2187$, and each digit can either be $1$ or $2$. So $2^1$ one digit numbers and so on up to $2^7$ $7$ digit.

 
Now we have to subtract out numbers from $2018$ to $2187$

 Then either the number must begin $221...$ or $222...$ with  four more digits at the end

 Using $1$s and $2$s there are $2^4$ options for each so:

 $2+4+8+16+32+64+128-2*16=256-2-32=\boxed{222}$

 





\textbf{Solution 3 (Casework)}

Since the greatest power of $3$ that can be used is $3^6$, we can do these cases.

 Coefficient of $3^6=0$: Then if the number has only $3^0$, it has 2 choices (1 or 2). Likewise if the number has both a $3^1$ and $3^0$ term, there are 4 choices, and so on until $3^5$, so the sum is $2+4+...+64=127-1=126$.

 Coefficient of $3^6=1$: Any combination of $1$ or $2$ for the remaining coefficients works, so $2^6=64$. Why is this less than $126$? Because the first time around, leading zeroes didn't count. But this time, all coefficients $3^n$ of $n<6$ need 1 and 2.

 Coefficient of $3^6=2$: Look at $3^5$ coefficient. If 1, all of them work because $3^7=2187-3^5=243<<2017$. That's 32 cases. Now of this coefficient is 2, then at the coefficient of $3^4=81$ is at least 1. However, $3^6*2+3^5*2+3^4>>2017$, so our answer is $126+64+32=\boxed{222}$.

 




	\item (\textit{AIME II 2017}) A set contains four numbers. The six pairwise sums of distinct elements of the set, in no particular order, are $189$, $320$, $287$, $234$, $x$, and $y$. Find the greatest possible value of $x+y$.


 
 \textbf{Answer: }791



\textbf{Solution 1}

Let these four numbers be $a$, $b$, $c$, and $d$, where $a>b>c>d$. $x+y$ needs to be maximized, so let $x=a+b$ and $y=a+c$ because these are the two largest pairwise sums. Now $x+y=2a+b+c$ needs to be maximized. Notice that $2a+b+c=3(a+b+c+d)-(a+2b+2c+3d)=3((a+c)+(b+d))-((a+d)+(b+c)+(b+d)+(c+d))$. No matter how the numbers $189$, $320$, $287$, and $234$ are assigned to the values $a+d$, $b+c$, $b+d$, and $c+d$, the sum $(a+d)+(b+c)+(b+d)+(c+d)$ will always be $189+320+287+234$. Therefore we need to maximize

 




	\item (\textit{AIME II 2017}) Find the sum of all positive integers $n$ such that $\sqrt{n^2+85n+2017}$ is an integer.


 
 \textbf{Answer: }195



\textbf{Solution 1}

Manipulating the given expression, $\sqrt{n^2+85n+2017}=\frac{1}{2}\sqrt{4n^2+340n+8068}=\frac{1}{2}\sqrt{(2n+85)^2+843}$. The expression under the radical must be an square number for the entire expression to be an integer, so $(2n+85)^2+843=s^2$. Rearranging, $s^2-(2n+85)^2=843$. By difference of squares, $(s-(2n+85))(s+(2n+85))=1\times843=3\times281$. It is easy to check that those are all the factor pairs of 843. Considering each factor pair separately, $2n+85$ is found to be $421$ and $139$. The two values of $n$ that satisfy one of the equations are $168$ and $27$. Summing these together gives us the answer ; $168+27=\boxed{195}$.

 





\textbf{Solution 2}

Clearly, the result when $n$ is plugged into the given expression is larger than $n$ itself. Let $x$ be the positive difference between that result and $n$, so that $\sqrt{n^2+85n+2017}=n+x$. Squaring both sides and canceling the $n^2$ terms gives $85n+2017=2xn+x^2$. Combining like terms, $(85-2x)n=x^2-2017$, so

 
 \[n=\frac{x^2-2017}{85-2x}.\]

 Since $n$ is positive, there are two cases, which are simple (luckily). Remembering that $x$ is a positive integer, then $x^2-2017$ and $85-2x$ are either both positive or both negative. The smallest value for which $x^2>2017$ is 45, which makes the denominator, and the entire expression, negative. Evaluating the other case where numerator and denominator are both negative, then we have that $x<45$ (from the numerator) and $85-2x<0$, which means $x>42$. This only gives two solutions, $x=43, 44$. Plugging these into the expression for $n$, we find that they result in 27 and 168, which both satisfy the initial question. Therefore, the answer is $168+27=\boxed{195}$.

 





\textbf{Solution 3 (Abuse the discriminant)}

Let the integer given by the square root be represented by $x$. Then $0 = n^2 + 85n + 2017 - x^2$. For this to have rational solutions for $n$ (checking whether they are integers is done later), the discriminant of this quadratic must be a perfect square. (This can be easily shown using the quadratic formula.)

 Thus, $b^2 - 4ac = 7225 + 4x^2 - 8068 = y^2$ for some integer $y$. Then $4x^2 - 843 = y^2$. Rearranging this equation yields that $843 = (2x+y)(2x-y)$. Noticing that there are 2 factor pairs of $843$, namely, $1*843$ and $3*281$, there are 2 systems to solve for $x$ and $y$ that create rational $n$. These yield solutions $(x,y)$ of $(211, 421)$ and $(71, 139)$.

 The solution to the initial quadratic in $n$ must then be $\frac{-85 \pm \sqrt{85^2 - 4(2017 - x^2)}}{2}$. Noticing that for each value of $x$ that has rational solutions for $n$, the corresponding value of the square root of the discriminant is $y$, the formula can be rewritten as $n = \frac{-85 \pm y}{2}$. One solution is $\frac{421 - 85}{2} = 168$ and the other solution is $\frac{139 - 85}{2} = 27$. Thus the answer is $168 + 27 = \boxed{195}$ as both rational solutions are integers.

 





\textbf{Solution 4 (Squeezing/Sandwich method)}

Notice that $(n+42)^2= n^2+84n+1764$. Also note that $(n+45)^2= n^2+90n+2025$. Thus, 
 \[(n+42)^2< n^2+85n+2017<(n+45)^2\] where $n^2+85n+2017$ is a perfect square. Hence,
 \[n^2+85n+2017= (n+43)^2\] or 
 \[n^2+85n+2017= (n+44)^2.\] Solving the two equations yields the two solutions $n= 168, 27$. Therefore, our answer is $\boxed{195}$.

 





\textbf{Solution 5 (Using factors)}

Let the expression be equal to $a$. This expression can be factored into $\sqrt{(n+44)^2-3n+81}$. Then square both sides, and the expression becomes $(n+44)^2-3n+81=a^2$. We have a difference of two squares. Rearranging, we have $(n+44+a)(n+44-a)=3(n-27)$. By inspection, the only possible values for $(n+44-a)$ are 0 and 1. When $(n+44-a)=0$, we must have $n-27=0$. Therefore, $27$ is a solution. When we have $(n+44-a)=1$, so $n=a-43$. Plugging this back to $(n+44+a)=3(n-27)$ (since $(n+44-a)=1$), we find that $a=211 \implies n=168$. Thus, the answer is $27+168= \boxed{195}$.

 \textbf{-RootThreeOverTwo}

 





\textbf{Solution 6}

Ignore the square root for now. This expression can be factored into $(n+44)^2-3n+81$. Just by inspection, when $n=27$, the expression becomes $71^2$, so $27$ is a solution. Proceed as Solution 5 to find the other solution(s).

 





\textbf{Solution 7 (alternative factoring)}

More intuitive, but a little bit slower considering the decimals.

 
Label the entire given expression as $k^2$.

 Instinctively we can do a crude completion of the square, resulting in $k^2$ = $(n+42.5)^2+210.75$Rearrange the equation to get a difference of squares.

 $k^2-(n+42.5)^2 = 210.75$

 $(k+n+42.5)(k-n-42.5) = 210.75$

 Factor $21,075$ to get $3^1$,$5^2$, and $281^1$

 Now the two factors given are either divided by 10 each or one being divided by 100. Let's start with the former case.

 If you try $281*3/10$ and $5*5/10$, you quickly realize that $n$ becomes negative. Naturally, you will realize you want the number's difference to be larger. Try $281*5/10$ and $3*5/10$. This gives an answer of $27$ for $x$. The next largest possibility also works, giving an $n$ of $168$. As you rise, some numbers don't work because it results in an n that is not an integer, as in the example of $\frac{281*5*5}{10}$ and $\frac{3}{10}$. 

 Now if you continue on with the next case, where one factor is divided by $100$, very swiftly will you realize most don't work simply because the difference is too small, or it doesn't give an integer. It helps a lot when you realize that the decimal does not end in a $5$, the answer will not be an integer. After a few short tests, we get $168+27=\boxed{195}$.

 


 -jackshi2006- minor latex edits by jske25

 




	\item (\textit{AIME II 2017}) Find the number of integer values of $k$ in the closed interval $[-500,500]$ for which the equation $\log(kx)=2\log(x+2)$ has exactly one real solution.


 
 \textbf{Answer: }501



\textbf{Solution 1}


                    \begin{figure}[H]
                        \centering
                        \includegraphics[width=0.5\textwidth]{Solution/AIME_2017_II_7_Solution_Figure_1.png}
                    \end{figure}
                    

 Note the equation $\log(kx)=2\log(x+2)$ is valid for $kx>0$ and $x>-2$. 

 Now, we simplify the given expression:$\log(kx)=2\log(x+2)=\log((x+2)^2) \implies kx=(x+2)^2$. The last equation is derived by taking away the outside logs from the previous equation. 

 Because $(x+2)^2$ is always non-negative, $kx$ must also be non-negative; therefore this takes care of the $kx>0$ condition as long as $k\neq0$, i.e. $k$ cannot be $0$. Now, we graph both $(x+2)^2$ (the green graph) and $kx$ (the red graph for $k=-1,k=-2,k=-3,k=8$) for $x>-2$. It is easy to see that all negative values of $k$ make the equation $\log(kx)=2\log(x+2)$ have only one solution. 

 However, there is also one positive value of $k$ that makes the equation only have one solution, as shown by the steepest line in the diagram. We can show that the slope of this line is a positive integer by setting the discriminant of the equation $(x+2)^2=kx$ to be $0$ and solving for $k$, which yields $k=8$. 

 When $1\leq k \leq 7$, $y=kx$ will intersect $y(x+2)^2$ twice, and when $k \geq 9$, $y=kx$ will never intersect $y(x+2)^2$.

 Therefore, there are $500$ negative solutions and $1$ positive solution, for a total of $\boxed{501}$.

 





\textbf{Solution 2}

We use an algebraic approach. Since $\log(kx)=2\log(x+2)$, then $kx = (x+2)^2$ (the converse isn't necessarily true!), or $x^2+(4-k)x+4=0$. Our original equation has exactly one solution if and only if there is only one solution to the above equation, or one of the solutions is extraneous; it involves the computation of the log of a nonpositive number.

 For the first case, this can only occur when it is a perfect square trinomial, or $k = 0, 8$. However, $k = 0$ results in $\log(0)$ on the left side, which is invalid. $k = 8$ yields $x = 2$, so that is one solution.

 For the second case, we can use the quadratic formula. We have 
 \[x = \frac{k-4 \pm \sqrt{k^2-8k}}2,\] so for there to be at least one real solution, the discriminant must be nonnegative, or $k < 0$ or $k > 8$. 

 Note that if $k > 8$, then both solutions to $x$ will be positive, and therefore both valid, which means all such $k$ are unsatisfactory. 

 We now wish to show that if $k < 0$, then exactly one solution works. Note that whenever $k < 0$, both "solutions" in $x$ are negative. One of the solutions to the equation is $x = \frac{k-4 + \sqrt{k^2-8k}}2$. We wish to prove that $x + 2 > 0$, or $x > -2$ (therefore the RHS in the original equation will be defined). Substituting, we have $\frac{k-4 + \sqrt{k^2-8k}}2 > -2$, or $\sqrt{k^2 - 8k} > -k$. Since both sides are positive, we can square both sides (if $k < 0$, then $-k > 0$) to get $k^2-8k > k^2$, or $8k < 0 \implies k < 0$, which was our original assumption, so this solution satisfies the original equation. The other case is when $x = \frac{k-4 - \sqrt{k^2-8k}}2$, which we wish to show is less that $-2$, or $\frac{k-4 - \sqrt{k^2-8k}}2 < -2 \implies k < \sqrt{k^2-8k}$. However, since the square root is defined to be positive, then this is always true, which implies that whenever $k < 0$, there is exactly one real solution that satisfies the original equation. Combining this with $k \in [-500, 500]$, we find that the answer is $500 + 1 = \boxed{501}$.

 \textbf{Note} The key to this solution is understanding that $\log(x+2)$ has a domain of $(-2, \infty),$ so in the second case, when there are two possible solutions of $x$ to $k<0,$ we notice that only the interval of the greater solution $(0,-2)$ works, which means that for any $k<0$ will have exactly one solution.

 \textbf{Reworded Solution 2} Immediately we notice if $k$ is non-zero we must have $kx, (x+2) > 0$ for our sole solution $x = x_0$. Simplifying the logarithmic equation we get $kx = (x+2)^2 \rightarrow 0 = x^2 + (4-k)x + 4 \rightarrow x = \frac{k-4 \pm \sqrt{k^2 - 8k}}{2}$. Then $k \leq 0$ or $k \geq 8$. When $k = 8$ we have exactly one real solution (easily verifiable). Notice when $k > 8$, both solutions to $x$ are positive, and so all such $k$ is not satisfactory. When $k < 0$ it can be shown that the greater solution to $x$ is in the interval $(0,-2)$ and the lesser solution is in the interval $(-2,-\infty)$ which is satisfactory. Then $k = 8,-1,-2,...,-500 \rightarrow \boxed{501}$ satisfactory integer values of $k$.

 




	\item (\textit{AIME II 2017}) Find the number of positive integers $n$ less than $2017$ such that 
 \[1+n+\frac{n^2}{2!}+\frac{n^3}{3!}+\frac{n^4}{4!}+\frac{n^5}{5!}+\frac{n^6}{6!}\] is an integer.


 
 \textbf{Answer: }134



\textbf{Solution 1}

We start with the last two terms of the polynomial $1+n+\frac{n^2}{2!}+\frac{n^3}{3!}+\frac{n^4}{4!}+\frac{n^5}{5!}+\frac{n^6}{6!}$, which are $\frac{n^5}{5!}+\frac{n^6}{6!}$. This can simplify to $\frac{6n^5+n^6}{720}$, which can further simplify to $\frac{n^5(6+n)}{720}$. Notice that the prime factorization of $720$ is $5\cdot3\cdot3\cdot2\cdot2\cdot2\cdot2$. In order for $\frac{n^5(6+n)}{720}$ to be an integer, one of the parts must divide $5, 3$, and $2$. Thus, one of the parts must be a multiple of $5, 3$, and $2$, and the LCM  of these three numbers is $30$. This means 
 \[n^5 \equiv 0\pmod{30}\] or 
 \[6+n \equiv 0\pmod{30}\] Thus, we can see that $n$ must equal $0\pmod{30}$ or $-6\pmod{30}$. Note that as long as we satisfy $\frac{6n^5+n^6}{720}$, $2!, 3!$, and $4!$ will all divide into integers, as their prime factorizations will be fulfilled with the LCM being 30. E.g. $4! = 2\cdot2\cdot2\cdot2\cdot3$, and this will be divisible by $2^4\cdot3^4\cdot5^4$. Now, since we know that $n$ must equal $0\pmod{30}$ or $-6\pmod{30}$ in order for the polynomial to be an integer, $n\equiv0, 24\pmod{30}$. To find how many integers fulfill the equation and are $<2017$, we take $\left \lfloor\frac{2017}{30} \right \rfloor$ and multiply it by $2$. Thus, we get $67\cdot2=\boxed{134}$.

 





\textbf{Solution 2}

Taking out the $1+n$ part of the expression and writing the remaining terms under a common denominator, we get $\frac{1}{720}(n^6+6n^5+30n^4+120n^3+360n^2)$. Therefore the expression $n^6+6n^5+30n^4+120n^3+360n^2$ must equal $720m$ for some positive integer $m$.Taking both sides mod $2$, the result is $n^6 \equiv 0 \pmod{2}$. Therefore $n$ must be even. If $n$ is even, that means $n$ can be written in the form $2a$ where $a$ is a positive integer. Replacing $n$ with $2a$ in the expression, $64a^6+192a^5+480a^4+960a^3+1440a^2$ is divisible by $16$ because each coefficient is divisible by $16$. Therefore, if $n$ is even, $n^6+6n^5+30n^4+120n^3+360n^2$ is divisible by $16$.

 Taking the equation $n^6+6n^5+30n^4+120n^3+360n^2=720m$ mod $3$, the result is $n^6 \equiv 0 \pmod{3}$. Therefore $n$ must be a multiple of $3$. If $n$ is a multiple of three, that means $n$ can be written in the form $3b$ where $b$ is a positive integer. Replacing $n$ with $3b$ in the expression, $729b^6+1458b^5+2430b^4+3240b^3+3240b^2$ is divisible by $9$ because each coefficient is divisible by $9$. Therefore, if $n$ is a multiple of $3$, $n^6+6n^5+30n^4+120n^3+360n^2$ is divisibly by $9$.

 Taking the equation $n^6+6n^5+30n^4+120n^3+360n^2=720m$ mod $5$, the result is $n^6+n^5 \equiv 0 \pmod{5}$. The only values of $n (\text{mod }5)$ that satisfy the equation are $n\equiv0(\text{mod }5)$ and $n\equiv4(\text{mod }5)$. Therefore if $n$ is $0$ or $4$ mod $5$, $n^6+6n^5+30n^4+120n^3+360n^2$ will be a multiple of $5$.

 The only way to get the expression $n^6+6n^5+30n^4+120n^3+360n^2$ to be divisible by $720=16 \cdot 9 \cdot 5$ is to have $n \equiv 0 \pmod{2}$, $n \equiv 0 \pmod{3}$, and $n \equiv 0 \text{ or } 4 \pmod{5}$. By the Chinese Remainder Theorem or simple guessing and checking, we see $n\equiv0,24 \pmod{30}$. Because no numbers between $2011$ and $2017$ are equivalent to $0$ or $24$ mod $30$ ($24+30\cdot0$ is the smallest and $0+30\cdot67$ is the largest), the answer is $\frac{2010}{30}\times2=\boxed{134}$.

 





\textbf{Solution 3}

Note that $1+n+\frac{n^2}{2!}+\frac{n^3}{3!}+\frac{n^4}{4!}+\frac{n^5}{5!}$ will have a denominator that divides $5!$. Therefore, for the expression to be an integer, $\frac{n^6}{6!}$ must have a denominator that divides $5!$. Thus, $6\mid n^6$, and $6\mid n$. Let $n=6m$. Substituting gives $1+6m+\frac{6^2m^2}{2!}+\frac{6^3m^3}{3!}+\frac{6^4m^4}{4!}+\frac{6^5m^5}{5!}+\frac{6^6m^6}{6!}$. Note that the first $5$ terms are integers, so it suffices for $\frac{6^5m^5}{5!}+\frac{6^6m^6}{6!}$ to be an integer. This simplifies to $\frac{6^5}{5!}m^5(m+1)=\frac{324}{5}m^5(m+1)$. It follows that $5\mid m^5(m+1)$. Therefore, $m$ is either $0$ or $4$ modulo $5$. However, we seek the number of $n$, and $n=6m$. By CRT, $n$ is either $0$ or $24$ modulo $30$, and the answer is $67+67=\boxed{134}$.

 -TheUltimate123

 





\textbf{Solution 4 (Step Solution)}

Clearly $1+n$ is an integer. The part we need to verify as an integer is, upon common denominator, $\frac{360n^2+120n^3+30n^4+6n^5+n^6}{720}$. Clearly, the numerator must be even for the fraction to be an integer. Therefore, $n^6$ is even and n is even, aka $n=2k$ for some integer $k$. Then, we can substitute $n=2k$ and see that $\frac{n^2}{2}$ is trivially integral. Then, substitute for the rest of the non-confirmed-integral terms and get $\frac{60k^3+30k^4+12k^5+4k^6}{45}$. It is also clear that for this to be an integer, which it needs to be, the numerator has to be divisible by 3. The only term we worry about is the $4k^6$, and we see that $k=3b$ for some integer $b$. From there we now know that $n=6b$. If we substitute again, we see that all parts except the last two fractions are trivially integral. In order for the last two fractions to sum to an integer we see that $n^5(6n+1) \equiv 0 \pmod5$, so combining with divisibility by 6, $n$ is $24$ or $0\pmod{30}$. There are $67$ cases for each, hence the answer $\boxed{134}$.

 




	\item (\textit{AIME II 2017}) A special deck of cards contains $49$ cards, each labeled with a number from $1$ to $7$ and colored with one of seven colors. Each number-color combination appears on exactly one card. Sharon will select a set of eight cards from the deck at random. Given that she gets at least one card of each color and at least one card with each number, the probability that Sharon can discard one of her cards and $\textit{still}$ have at least one card of each color and at least one card with each number is $\frac{p}{q}$, where $p$ and $q$ are relatively prime positive integers. Find $p+q$.


 
 \textbf{Answer: }013



\textbf{Solution 1}

Without loss of generality, assume that the $8$ numbers on Sharon's cards are $1$, $1$, $2$, $3$, $4$, $5$, $6$, and $7$, in that order, and assume the $8$ colors are red, red, and six different arbitrary colors. There are ${8\choose2}-1$ ways of assigning the two red cards to the $8$ numbers; we subtract $1$ because we cannot assign the two reds to the two $1$'s.In order for Sharon to be able to remove at least one card and still have at least one card of each color, one of the reds have to be assigned with one of the $1$s. The number of ways for this to happen is $2 \cdot 6 = 12$ (the first card she draws has to be a $1$ (2 choices), while the second card can be any card but the remaining card with a $1$ (6 choices)). Each of these assignments is equally likely, so desired probability is $\frac{12}{{8\choose2}-1}=\frac{4}{9} \implies 4 + 9 = 13 = \boxed{013}$.

 





\textbf{Solution 2}

First note that out of the $8$ selected cards, one pair of cards have to share the same number and another pair of cards have to share the same color. Now, these $2$ pairs of cards can't be the same or else there will be $2$ cards which are completely same. Then, WLOG let the numbers be $1,1,2,3,4,5,6,$ and $7$ and the colors be $a,a,b,c,d,e,f,$ and $g$. We therefore obtain only $2$ cases:

 Case One: $1a,1b,2a,3c,4d,5e,6f,$ and $7g$In this case, we can discard $1a$.There are $2*6=12$ situations in this case.

 Case Two: $1b,1c,2a,3a,4d,5e,6f,$ and $7g$In this case, we can't discard.There are $\dbinom{6}{2}=15$ situations in this case.

 So the probability is $\frac{12}{12+15}=\frac{4}{9}$, giving us the answer of $4+9=\boxed{013}$.

 





\textbf{Solution 3}

There are $7!$ ways to choose a set of 7 cards that have all the numbers from 1-7 and all 7 colors. There are then $42$ cards remaining. Thus, there are $7!(42)$ desired sets.

 Now, the next thing to find is the number of ways to choose 8 cards where there is not a set of 7 such cards. In this case, one color must have 2 cards and one number must have 2 cards, and they can't be the same number/color card. The number of ways to pick this is equal to a multiplication of $\binom{7}{2}$ ways to pick 2 numbers, $7$ colors to assign them to, $\binom{6}{2}$ ways to pick 2 nonchosen colors, $5$ ways to pick a number to assign them to, and $4!$ ways to assign the rest.

 Thus, the answer is $\frac{7!(42)}{7!(42) + 21(15)(7)(5!)}$. Dividing out $5!$ yields $\frac{42(42)}{42(42) + 21(15)(7)}$ which is equal to $\frac{2(42)}{2(42) + 15(7)}$ which is equal to $\frac{12}{12 + 15}$ which is equal to $\frac{4}{9}$ giving a final answer of $\boxed{013}$.

 





\textbf{Solution 4}

We can rewrite the problem as "What is the probability that, given 8 cards with numbers a, b, c, d, e, f, g, g and some assortment of seven colors A, B, C, D, E, F, G, G where G is repeated, one of the cards is Gg?" Note that this is a valid restatement because Sharon has to have two of one number and two of one color. She needs to be able to take away one of the cards with the duplicate number, but this also has to have the duplicate color. There are two cases.

 Case I: One of the cards is Gg. This implies that the other card with color G can be placed in $6$ ways, and the rest of the colors can be paired with cards in $6!$ ways. 

 Case II: None of the cards are Gg. This implies that the cards with color G can be chosen in $\dbinom{6}{2}=15$ ways, and the rest of the colors can be paired with cards in $\frac{6!}{2}$ ways, with the divide by 2 because of the double-g.

 Note that there is no case with Gg, Gg because all 49 cards are unique!

 Therefore our answer is $\frac{6\times6!}{\frac{15}{2}\times6! + 6\times6!}=\frac{4}{9}$, so $\boxed{013}$.

 





\textbf{Solution 5}

We count the two entities: The number of sets of eight cards that contain all seven numbers and all seven colors; and The number of sets of eight cards that contain all seven numbers and colors, and one card can be removed and the property still holds.

 For the first: It is equivalent to counting:How many ways can (colors) $A, A, B, C, D, E, F, G$ be matched to (numbers) $1, 1, 2, 3, 4, 5, 6, 7$ such that the two $A's$ cannot both be matched to the two ones. (This would mean we've chosen two identical cards from the deck, both with color $A$ and number one.)We consider two cases: 

 1) Neither $A$ is matched with a one. Then there are $\binom{6}{2}$ ways to choose which two numbers are matched with the $A's,$ and $6! / 2!$ ways to order the remaining numbers. (There are two one's, so we divide by two.)

 2) One of the $A$'s is matched with a one. There are $6 \cdot 6!$ ways to arrange the remaining color

 Since there are seven ways to choose the extra color and likewise for the extra number, we have in total $7 \cdot 7 \cdot \left( 6C2 \cdot \frac{6!}{2!} + 6 \cdot 6! \right)$ ways to choose eight cards satisfying the first condition.

 For the second: View the cards as rooks on a seven by seven chess board. Seven cards that contain all colors and numbers is can be represented as placing seven rooks on the board so that no two rooks are attacking each other. There are $7!$ ways to do this. For each arrangement, we can choose any of the other 42 positions on the board (cards) to add an extra rook. Thus, there are $7! \cdot 42$ good sets of eight cards. Note that we are not overcounting.

 Thus, the probability of getting a good set given the first condition is $\frac{ 7! \cdot 42  }{ 7 \cdot 7 \cdot \left( 6C2 \cdot \frac{6!}{2!} + 6 \cdot 6! \right) } = \frac{4}{9},$ and the answer is $\boxed{013}.$

 





\textbf{Solution 6}

We recast this problem to coloring 8 cells in a 7x7 grid, such that each row and column has at least one colored cell. If a cell in row a and column b is colored, that means we drew a card with color a and number b. 

 By pigeonhole principle, there exists a row that has 2 colored cells, and intuitively, that row is unique. There also exists a column with 2 colored cells. Call the unique row, "Row A", and Call the unique column, "Column B".

 If we are able to remove a card so that we still have a card of each color and number, that means, we can erase a cell, so each row still has at least one colored cell, and so does each column. If we erase a cell not in Row A, then the erased cell's row will be empty. If we erase a cell not in Column B, the erased cell's column will be empty. Thus the cell that we erase must be in both Row A and Column B. There is only one cell in both Row A and Column B, and thus that cell must be colored.

 We first count the number of successful outcomes. We must chose which row to be "Row A", and which column to be "Column B". There are 7 ways to chose "Row A", and 7 ways to chose "Column B". Because of the logic used above, the intersection of Row A and Column B must be colored. However, Row A needs to have 2 colored cells. There are 6 other uncolored cells to chose from. Similarly, there are 6 uncolored cells in Column B. 

 Note that we can ignore all the rows and columns that contain colored cells. This is because Row A already contains 2 colored cells, and thus we can't color any more cells in Row A, similarly for Column B. The other row that contains a colored cell is not Row A, so it can only contain 1 colored cell, which it has, thus we can't color anymore cells in it. Similarly for the other column.  

 We can actually remove these rows/columns. This leaves a 5x5 grid, full of uncolored cells. Each row can have 1 cell, and so can each column. We count the \# of colorings from the perspective of the rows. We look at the first row, it has 5 choices. We can look at the 2nd row, it has 4 choices, as it's selected cell can't occupy the same column. We continue down the line, and see there are $5! = 120$ ways to color.

 Thus, the \# of successful outcomes is $5! \cdot 6 \cdot 6 \cdot 7 \cdot 7$.

 We can then count the \# of possible outcomes. 

 We can count the \# of unsuccessful outcomes. We select "Row A" and "Column B" as usual, with $7 \cdot 7$ ways. We don't color the intersection, so we have to color 2 cells in each. There are 6 other cells that can be colored in each, so we have ${6 \choose 2} = 15$ ways to color each. Then, we can get rid of rows and columns using the method above and we are left with a 4x4 grid. There are $4! = 24$ ways to color. 

 Thus, the \# of unsuccessful outcomes is $4! \cdot 15 \cdot 15 \cdot 7 \cdot 7$.

 After some calculation, we get $\frac{successful}{possible} = \frac{successful}{successful+unsuccessful} = \frac{5! \cdot 6^2 \cdot 7^2}{7^2 \cdot(5! \cdot 6^2 + 4! \cdot 15^2} = \frac{5! \cdot 6^2}{5! \cdot 6^2 +4! \cdot 15^2} = \frac{5 \cdot 6^2}{5 \cdot 6^2 + 15^2} = \frac{180}{180+225} = \frac{4}{9} \implies 4+9 = \boxed{13}$

 -Alexlikemath

 





\textbf{Solution 7 (Official MAA)}

In a set of eight cards that includes every color and every number, there will be exactly one repeated number and exactly one repeated color. If Sharon selects a set that includes the card with that number and color, Sharon can discard it. If the set does not include that card, Sharon cannot discard any card.In the former case, there are $7!$ ways to choose one card of each number and color, and then $42$ ways to choose the "extra" card. The total is $7!\cdot 42$.In the latter case, there are $7$ ways to choose which number is repeated, and $\binom{7}{2} = 21$ ways to choose which two cards of that number are used. There are then $5$ ways to choose which color is repeated, and $\binom{6}{2} = 15$ ways to choose which cards of that color are used. There now remain four unused colors and four unused numbers, and there are $4! = 24$ ways to choose from those. The total is $7 \cdot 21 \cdot 5 \cdot 15 \cdot 24$.The probability that a given set is of the first type is therefore $\frac{7! \cdot 42}{7! \cdot 42+7 \cdot 21 \cdot 5 \cdot 15 \cdot 24} = \frac{4}{9}.$ The requested sum is $4 + 9 = 13$.

 





\textbf{Solution 8}

WLOG let the colors be $R, R, O, Y, G, B, I, V$ and let the numbers be $1, 1, 2, 3, 4, 5, 6, 7$. Now call a color or number $safe$ is there is at least two cards in that category. In this case, the color red is $safe$ and the number $1$ is safe because we have two of each. 

 Now notice that a card can be removed if and only if both its color AND number are $safe$, since even if we remove it we will have one of each category. Therefore, we must have a $R1$. 

 
For this probability, we will first count the total ways to assign the cards. Treating each color and number as distinct, we have $8!$ ways. However, we cannot have two red ones, so we must subtract $2 * 6!$ cases from our total. 

 Meanwhile, to count the successful outcomes, one of the ones must be red. We have $2$ choices for the one, and $2$ choices for the red. Then, the other one has $6$ colors, and there are $6!$ ways to color the rest of the numbers.

 
Diving and simplifying, we have $\frac{2*2*6*6!}{8!-2*6!} = \frac{4}{9} \implies 4+9 = \boxed{13}$.

 




	\item (\textit{AIME II 2017}) Rectangle $ABCD$ has side lengths $AB=84$ and $AD=42$. Point $M$ is the midpoint of $\overline{AD}$, point $N$ is the trisection point of $\overline{AB}$ closer to $A$, and point $O$ is the intersection of $\overline{CM}$ and $\overline{DN}$. Point $P$ lies on the quadrilateral $BCON$, and $\overline{BP}$ bisects the area of $BCON$. Find the area of $\triangle CDP$.


 
 \textbf{Answer: }546



\textbf{Solution 1}


                    \begin{figure}[H]
                        \centering
                        \includegraphics[width=0.5\textwidth]{Solution/AIME_2017_II_10_Solution_Figure_1.png}
                    \end{figure}
                    Impose a coordinate system on the diagram where point $D$ is the origin. Therefore $A=(0,42)$, $B=(84,42)$, $C=(84,0)$, and $D=(0,0)$. Because $M$ is a midpoint and $N$ is a trisection point, $M=(0,21)$ and $N=(28,42)$. The equation for line $DN$ is $y=\frac{3}{2}x$ and the equation for line $CM$ is $\frac{1}{84}x+\frac{1}{21}y=1$, so their intersection, point $O$, is $(12,18)$. Using the shoelace formula on quadrilateral $BCON$, or drawing diagonal $\overline{BO}$ and using $\frac12bh$, we find that its area is $2184$. Therefore the area of triangle $BCP$ is $\frac{2184}{2} = 1092$. Using $A = \frac 12 bh$, we get $2184 = 42h$. Simplifying, we get $h = 52$. This means that the x-coordinate of $P = 84 - 52 = 32$. Since P lies on $\frac{1}{84}x+\frac{1}{21}y=1$, you can solve and get that the y-coordinate of $P$ is $13$. Therefore the area of $CDP$ is $\frac{1}{2} \cdot 13 \cdot 84=\boxed{546}$.

 





\textbf{Solution 2 (No Coordinates)}

Since the problem tells us that segment $\overline{BP}$ bisects the area of quadrilateral $BCON$, let us compute the area of $BCON$ by subtracting the areas of $\triangle{AND}$ and $\triangle{DOC}$ from rectangle $ABCD$.

 To do this, drop altitude $\overline{OE}$ onto side $\overline{DC}$ and draw a segment $\overline{MQ}$ parallel to $\overline{AN}$ from side $\overline{AD}$ to $\overline{ND}$. Since $M$ is the midpoint of side $\overline{AD}$, 
 \[\overline{MQ}=14\] Denote $\overline{OE}$ as $a$. Noting that $\triangle{MOQ}~\triangle{COD}$, we can write the statement 
 \[\frac{\overline{DC}}{a}=\frac{\overline{MQ}}{21-a}\] 
 \[\implies \frac{84}{a}=\frac{14}{21-a}\] 
 \[\implies a=18\] Using this information, the area of $\triangle{DOC}$ and $\triangle{AND}$ are 
 \[\frac{18\cdot 84}{2}=756\] and 
 \[\frac{28\cdot 42}{2}=588\] respectively. Thus, the area of quadrilateral $BCON$ is 
 \[84\cdot 42-588-756=2184\] Now, it is clear that point $P$ lies on side $\overline{MC}$, so the area of $\triangle{BPC}$ is 
 \[\frac{2184}{2}=1092\] Given this, drop altitude $\overline{PF}$ (let's call it $b$) onto $\overline{BC}$. Therefore, 
 \[\frac{42b}{2}=1092\implies b=52\] From here, drop an altitude $\overline{PG}$ onto $\overline{DC}$. Recognizing that $\overline{PF}=\overline{GC}$ and that $\triangle{MDC}$ and $\triangle{PGC}$ are similar, we write 
 \[\frac{\overline{PG}}{\overline{GC}}=\frac{\overline{MD}}{\overline{DC}}\] 
 \[\implies \frac{\overline{PG}}{52}=\frac{21}{84}\] 
 \[\implies \overline{PG}=13\] The area of $\triangle{CDP}$ is given by 
 \[\frac{\overline{DC}\cdot \overline{PG}}{2}=\frac{84\cdot 13}{2}=\boxed{546}\] ~blitzkrieg21 and jdong2006

 





\textbf{Solution 3 (Coord bash)}

Let $D = (0,0), C = (84,0), B = (84,42), A = (0,42)$. We can also find that $M = (0,21)$ by the midpoint formula. Now, lets find $O$. $O$ is the intersection of $ND$ and $CM$. So, we just need to find the equation of those lines, and solve. For $DN$, we first find the slope to be $\dfrac{42}{28} = \dfrac{3}{2}$. And obviously, the y-int is just $0$. So, our equation is: 
 \[y=\dfrac{3}{2}x\]. Now, for $CM$, we have a slope of: 
 \[\dfrac{21}{-84} = -\dfrac{1}{4}\], and the y-int is just $21$, so our equation is: 
 \[y=-\dfrac{1}{4}x+21\], solving for $x$ and $y$, we get: $(12,18)$. Now, we use the fact that $\overline{BP}$ bisects the area of $BCON$. By shoelace, we find the area of $BCON$ to be $2184$. Now, we will use $\triangle{BPC}$. We have that its area is $\dfrac{2184}{2} = 1092$. So, by the area of a triangle formula the height $BCP$ is $52$. So, the $x$ coordinate of $P$ is $84-52 = 32$. Then, we use the fact that $P$ lies on $CM$, so using the equation of line $CM$ (which we found to be $y=-\dfrac{1}{4}x+21$, plugging $x=32$, we get $y = 13$, so we have the coordinates of $P$ to be $(32,13)$, so the height of $PCD$ is $13$, so the area is $\dfrac{13\cdot84}{2}=\boxed{546}$

 





\textbf{Solution 4 (Menelaus, Similarity)}


                    \begin{figure}[H]
                        \centering
                        \includegraphics[width=0.5\textwidth]{Solution/AIME_2017_II_10_Solution_Figure_2.png}
                    \end{figure}
                    From the givens, we know that $DM=MA=21$, $AN=28$, and $BN=56$. Extend $BA$ to meet $CM$ at point $X$. Since $M$ is the midpoint of $AD$, $\triangle{AXM}$ and $\triangle{DCM}$ are congruent, so $AX=84$. From Menelaus' Theorem on $\triangle{DAN}$, we have
 \[\dfrac{DM}{MA} \cdot \dfrac{AX}{NX}\cdot \dfrac{NO}{OD}=1.\]After plugging in the known values, we have that $\frac{NO}{OD}=\frac{4}{3}$. Let point $Q$ be the point on $DA$ such that $\triangle{DQO}\sim\triangle{DAN}$. Then, because of the ratio we just found, we know that $QO=12$. Thus, if we let point $Q^{\prime}$ be the point on $BC$ such that $OQ^{\prime} \parallel AB$, we have $OQ^{\prime}=72.$ This means that $[COB]=\frac{1}{2}(72)(42).$ We can do the same process vertically to find that $[NOB]=\frac{1}{2}(24)(56)$. Summing, we have that $[BCON]=12\cdot 182$. Then, $[CPB]=6\cdot 182$. Since $BC=42$, the altitude from $P$ to $BC$ is $52$. Let the altitude from $P$ to $CD$ be $h$. Then, from similar triangles,
 \[\dfrac{h}{21}=\dfrac{52}{84}.\]This gives that $h=13$, and thus the area is $\frac{1}{2}(84)(13)=\boxed{546}$.

 - The\_Other\_Guy

 




	\item (\textit{AIME II 2017}) Five towns are connected by a system of roads. There is exactly one road connecting each pair of towns. Find the number of ways there are to make all the roads one-way in such a way that it is still possible to get from any town to any other town using the roads (possibly passing through other towns on the way).


 
 \textbf{Answer: }544



\textbf{Solution}

It is obvious that any configuration of one-way roads which contains a town whose roads all lead into it or lead out of it cannot satisfy the given. We claim that any configuration which does not have a town whose roads all lead into it or lead out of it does satisfy the given conditions. Now we show that a loop of $3$ or more towns exist. Pick a town, then choose a neighboring town to travel to $5$ times. Of these $6$ towns visited, at least two of them must be the same; therefore there must exist a loop of $3$ or more towns because a loop of $2$ towns cannot exist. We want to show that the loop can be reached from any town, and any town can be reached from the loop.

 $\textbf{Case 1}$. The loop has $5$ towns.Clearly every town can be reached by going around the loop.

 $\textbf{Case 2}$. The loop has $4$ towns.The town not on the loop must have a road leading to it. This road comes from a town on the loop. Therefore this town can be reached from the loop. This town not on the loop must also have a road leading out of it. This road leads to a town on the loop. Therefore the loop can be reached from the town.

 $\textbf{Case 3}$. The loop has $3$ towns.There are two towns not on the loop; call them Town $A$ and Town $B$. Without loss of generality assume $A$ leads to $B$. Because a road must lead to $A$, the town where this road comes from must be on the loop. Therefore $A$ and therefore $B$ can be reached from the loop. Because a road must lead out of $B$, the town it leads to must be on the loop. Therefore the loop can be reached from $B$ and also $A$.

 The number of good configurations is the total number of configurations minus the number of bad configurations. There are $2^{{5\choose2}}$ total configurations. To find the number of bad configurations in which a town exists such that all roads lead to it, there are $5$ ways to choose this town and $2^6$ ways to assign the six other roads that do not connect to this town. The same logic is used to find the number of bad configurations in which a town exists such that all roads lead out of it. It might be tempting to conclude that there are $5 \cdot 2^6+5 \cdot 2^6$ bad configurations, but the configurations in which there exists a town such that all roads lead to it and a town such that all roads lead out of it are overcounted. There are $5$ ways to choose the town for which all roads lead to it, $4$ ways to choose the town for which all roads lead out of it, and $2^3$ ways to assign the remaining $3$ roads not connected to either of these towns. Therefore, the answer is $2^{{5\choose2}}-(5 \cdot 2^6+5 \cdot 2^6-5\cdot 4 \cdot 2^3)=\boxed{544}$.

 \textbf{Diagram} 
                    \begin{figure}[H]
                        \centering
                        \includegraphics[width=0.5\textwidth]{Solution/AIME_2017_II_11_Solution_Figure_1.JPG}
                    \end{figure}
                    

 Note that this isn't necessarily a diagram of the actual sequence of roads (I haven't drawn in every pair of cities to be connected), but a simplified diagram that can be used to represent a visualized representation of the loops in the cases\textbf{mathboy282}

 





\textbf{Solution 2 (complementary counting)}

The only way a town does not meet the conditions in the question is if the town has either all roads leading towards it or all roads leading away from it. For example, if all roads lead away from Town $A$, there is no way to reach the town starting from Towns $B$, $C$, $D$, or $E$. If all roads lead towards Town $A$, there is no way to reach any other town starting from Town $A$. Thus, we will count the ways this occurs. (To show this, first note that there must always be a directed path that visits every town once, say $A \rightarrow C \rightarrow B \rightarrow E \rightarrow D.$ (Redei's Theorem). Taking cases, it can be seen that if $A$ does not have all its roads leading, then $E$ must have its roads leading into it.)

 
$\textbf{Case 1}$. WLOG, Town $A$ has all roads leading towards it.

 In this case, the four roads leading to Town $A$ have already been determined. There are $6$ roads that have not been given directions. Each of these roads have $2$ options: it can lead towards one town or the other town. Thus, there are $2^6$ arrangements of roads.However, we must consider the $5$ towns in which this scenario can occur: there are $5 \cdot 2^6=320$ arrangements.

 
$\textbf{Case 2}$. WLOG, Town $A$ has all roads leading away.

 Notice that this case is symmetrical to the first case. Thus, there are $320$ arrangements.

 
$\textbf{Case 3}$. WLOG, Town $A$ has all roads leading towards it and Town $B$ has all roads leading away.

 We must check for double-counted cases. Drawing a flow diagram, we see that this case determines $7$ roads. For the undetermined $3$ roads, there are $2^3=8$ arrangements. However, we must again consider the different ways that this case can occur. There are $5$ choices for towns with roads leading towards it $4$ choices for roads leading away. Thus, the total double-counted cases are $5 \cdot 4 \cdot 8=160$ arrangements.

 We must subtract the cases we counted from the total. There are a total of $10$ roads with $2$ possible arrangements each.Therefore the total number of cases is $2^{10} =1024$, and the number of cases that meet the conditions is $1024 - (320+320-160) = \boxed{544}$ arrangements.

 





\textbf{Solution 3}

Assume the five towns are named A, B, C, D, and E. Draw roads that connect each of the five towns. First, figure out how many total roads there are with no restrictions. This would be $2^{10}$, or $1024$. Next, add together the number of restrictions (or incorrect cases) there are. One incorrect case happens if we can't get to one of the five towns. When that happens, any roads that connect that town with any of the other four towns must be one way. Four roads are invalidated (because they are one way) as, assuming we can't get to town A, roads AB, AC, AD, and AE would all be one way. Since there are ${{5\choose1}}$ ways of choosing which town to restrict, there are $5\cdot2^{10-4}$, or $320$ restrictions in this case. The next case is when we can't get to two of the towns. Let's assume we don't want to get to towns A and B. Then, roads AC, AD, AE, BC, BD, and BE are all invalidated.  Since there are ${{5\choose2}}$ of choosing which two towns to restrict, there are $10\cdot2^{10-6}$, or $160$ restrictions in this case. Notice that when we restrict three towns, it is the same thing as restricting two towns. For example, if we restricted entry to towns C, D, and E from towns A and B, it's the same thing as if we restricted entry to towns A and B from towns C, D, and E, as the exact same roads are invalidated in both cases. Thus, our final answer is $1024-(320+160) = \boxed{544}$ arrangements.

 Solution by IronicNinja~

 





\textbf{Solution 4}

As noted before, you can see that there are 2 ways that the condition cannot be met; it either has all roads leading into one city, or all roads leading out of one city. We then use complementary counting to count these cases, with PIE. Obviously there are $2^5*2^5$ ways to make roads (just draw a pentagon with all of the diagonals and assume that they are roads, and count the sides and then the diagonals). Then, make all roads leading in to another city. There are 4 roads leading into that city, so there are $2^6$ ways to designate that. Given that there are 5 cities, and it can also be the ways to go out of the city, you get $10*2^6$ ways NOT to get into the city. However, this disregards the case that 2 cities are like this, which is the PIE case. Quick inspection shows that there are no cases where 3 are like that, and inspection with le pentagon can also show you that there are $\binom{5}{2}*2$ (because they can all lead in and all lead out) $*2^3$ (because there are 3 roads left to choose). In the end, this comes out to $2^5*2^5-(10*2^6-2*\binom{5}{2}*2^3)=1024-(640-160)=\boxed{544}$

 -dragonon

 





\textbf{Solution 5 (Graph Theory)}

We use Claim 1 along with complementary counting.Then, the number of valid cases is the total ($2^{10}$) minus the cases in which there are towns with indegree or outdegree 5.Number of cases in which there is a node with indegree 5: $5*2^6$ 5 ways to choose which town has indegree 5 and for the remaining roads, they each have 2 possible statuses.Similarly, the number of cases in which there is a node with outdegree 5 is also $5*2^6$The number of cases in which there is a node with indegree 5 and there is a node with outdegree 5: $5*4*2^3$

 Thus, answer = $2^{10}-(5*2^6+5*2^6-5*4*2^3)$ = $544$

 
                \begin{itemize}
                    \item Note: In this graph there can be at most one town with indegree 5 or outdegree 5.
                \end{itemize}


                Solution by ZHANALEX

 \textbf{Claim 1} Claim: If and only if there is no town in this graph with a indegree or outdegree of 4, it is possible to reach all of the 4 other towns from any town in the graph.

 Proof:

 Part 1: If there is no town in this graph with a indegree or outdegree of 4, it is possible to reach all of the 4 other towns from any town in the graph

 Since every town has an outgoing road, there is definitely a cycle. (Since there are only n town, eventually you must come back to a town that is already visited)Since a cycle of length 2 is not possible in this problem, the cycle can have a length of 3, 4, or 5.If the cycle has length 5, then it is definitely possible to go from one town to any other town.Ex: If the cycle is 1->2->3->4->5->1, then it is clear that you can reach any other town from one town

 If the cycle has length 4:WLOG, let 1->2->3->4->1 be the cycle. Since town 5 has at least one incoming road and one outgoing road, there is definitely an road from a town on the cycle to town 5. There must also be an road from town 5 to another town that is in the cycle. Thus, it is definitely possible to reach any other town from one town.

 If the cycle has length 3:WLOG, we can assume that the cycle is 1->2->3->1 and that the road between town 4 and town 5 is 4->5. Since 4 has at least one incoming road, there must be a town in the cycle that has an road to town 4. Since town 5 has at least one outgoing road, there must be a road from town 5 to some town in the cycle. Thus, it is possible to reach any other town starting from one town.

 
Part 2: If it is possible to reach all of the 4 other towns starting from any town, there can be no town with indegree or outdegree of 4.It is clear that you cannot have any town with indegree or outdegree of 4 or else that town would either be unreachable from the other towns or that town would be unable to reach any other town.

 The proof is complete.(It is actually much simpler than it looks)

 -ZHANALEX

 





\textbf{Solution 6}

We claim the following conditions:

 -A path is invalid if a vertex $\omega$ has no arcs leading to it. (1)

 -A path is invalid if a vertex $\omega$ has all arcs leading to it. (2)

 -A path is always valid if it does not satisfy (1) and (2) and has a "cycle" of arcs of $\omega_1 \to \omega_2 \to \omega_3 \to \omega_1$. (3)

 -(1) and (2) are sufficient to cover all invalid paths.

 Call the vertices $\omega_1, \omega_2, \omega_3, \omega_4, \omega_5$. WLOG, let $\omega_1 \to \omega_2$. Then, WLOG let $\omega_2 \to \omega_3$. We have two ways to go from here:

 If $\omega_3 \to \omega_1$, we call it a "cycle".

 We will prove (3):

 Notice that once $\omega_4, \omega_5$ can reach one of the vertices on the "cycle", it can reach any of the other vertices on the same cycle simply by travelling along the arcs of the cycle. Additionally, one of $\omega_4, \omega_5$ must be able to travel to the cycle. Consider that it cannot, then $\omega_3 \to \omega_4$ and $\omega_5 \to \omega_4$ are forced. Now, $\omega_4 \to \omega_1, \omega_2, \omega_3$ which is on the cycle. Now, if $\omega_5 \to \omega_4$, it will also be able to travel to the cycle.  Otherwise, the directed arc will directly travel to the cycle. Also, each of the vertices $\omega_1, \omega_2, \omega_3$ will travel to $\omega_4, \omega_5$. 

 If our graph does not satisfy (1) and (2), $\omega_4, \omega_5$ will connect to the cycle on both directions. This can be shown below:

 Clearly $\omega_4, \omega_5$ are connected. Assuming the graph does not satisfy (1) and (2), $\omega_4, \omega_5$ are connected to the cycle consisting of $\omega_1, \omega_2, \omega_3$. WLOG, let $\omega_4 \to \omega_5$:

 Again, to avoid (2), $\omega_5 \to \omega_a$ where $a\in{1,2,3}$. Now to avoid (1), one of $\omega_a \to \omega_4, a\in{1,2,3}$.

 Now it is clear that each vertex can go to any other vertex. This concludes our proof.

 Say we try to avoid (3) so that WLOG $\omega_5 \to \omega_1$. If we keep trying to avoid a "cycle", $\omega_3 \to \omega_4$ and $\omega_4 \to \omega_5$ which is a cycle. 

 Therefore, if there is any path $\omega_a \to \omega_b \to \omega_c$, it will lead to a cycle and ultimately a valid path. The only way to avoid this is to have (1) or (2).

 Consider (1):

 WLOG, let $\omega_1$ have no arcs leading to it. Then all other vertices will have an arc directed toward them away from $\omega_1$. Therefore, one path can have only 1 unique vertex that has no arcs leading to it. There are 5 ways to choose the vertex and $2^{\binom{5}{2}-4}=2^6$ ways for the remaining 6 arcs since they can be directed in any way. There are $2^6\cdot5=320$ paths for this condition.

 Consider (2):

 WLOG, let $\omega_1$ have all arcs leading to it. Similarly, all other vertices will have an arc that is directed away from them and points to $\omega_1$. Therefore, only one unique vertex can have all arcs leading to it. There are 5 ways to choose the vertex and $2^{\binom{5}{2}-4}=2^6$ ways for the remaining arcs. There are also $2^6\cdot5=320$ paths for this condition.

 However, we have overcounted paths that satisfy both (1) and (2). These paths will have one vertex that has all arcs leading to it and another that has all arcs directed away from it. There are $5\cdot4=20$ ways to choose these 2 vertices, and the remaining $\binom{3}{2}=3$ arcs each have 2 ways to be directed toward each other, so there are a total of $2^3\cdot 20=160$ overlapping paths.

 By PIE, the number of invalid cases is $320+320-160=480$.

 The total number of cases is $2^{\binom{5}{2}}=2^{10}=1024$ since there are 10 edges and each can go 2 ways. By complementary counting, the number of valid paths is just $1024-480=\boxed{544}$

 -Magnetoninja

 





\textbf{Solution 7}

Let $\mathbf{\Omega}_n$ denote the $n$-th total direction tournament set. Since every mapping between two nodes is determined in $\binom{n}{2}$ ways, and two directions each, we have $2^{\binom{n}{2}}$ total possible cases. For $n=5, \mathbf{\Omega}_5 = 1024$. Define a \textbf{component} to be a cluster of isolated towns (so basically clustered sample for those fluent in AP stats haha, i feel bad for you...).

 Let $\pi_n$ denote the number of strict tournament sets. We wish to find $\pi_5$. Recall that $\mathbf{\Omega}_n$ can be chosen via picking a subset of $k$ towns to form the first component, and letting the remaining $n-k$ towns form any tournament. Thus,

 
 \[\mathbf{\Omega}_n = \sum_{k=1}^{n} \binom{n}{k} \pi_n \mathbf{\Omega}_{n-k}\].

 Note that we can never pick no node as at least one component must be chosen to be strict.

 Note that $\mathbf{\Omega}_0 = 1$, so we can just solve via the sum.

 
 \[\mathbf{\Omega}_1 = 1 \cdot \pi_1 \cdot 1 \implies 1 = \pi_1\]
 \[\mathbf{\Omega}_2 = 2 \cdot \pi_1 \cdot 1 + 2 \cdot \pi_2 \cdot 1 \implies 0 = \pi_2\]
 \[\mathbf{\Omega}_3 = 3 \cdot \pi_1 \cdot 2 + 3 \cdot \pi_2 \cdot 1 + 1 \cdot \pi_3 \cdot 1 \implies 2 = \pi_3\]
 \[\mathbf{\Omega}_4 = 4 \cdot \pi_1 \cdot 8 + 6 \cdot \pi_2 \cdot 2 + 4 \cdot \pi_3 \cdot 1 + 1 \cdot \pi_4 \cdot 1 \implies 24 = \pi_4\]
 \[\mathbf{\Omega}_ 5 = 5 \cdot \pi_1 \cdot 64 + 10 \cdot \pi_2 \cdot 8 + 10 \cdot \pi_3 \cdot 2 + 5 \cdot \pi_4 \cdot 1 + 5 \cdot \pi_5 \cdot 1 \implies 1024 = 480 + \pi_5\]

 Therefore, $\pi_5 = \boxed{544}$.

 




	\item (\textit{AIME II 2017}) Circle $C_0$ has radius $1$, and the point $A_0$ is a point on the circle. Circle $C_1$ has radius $r<1$ and is internally tangent to $C_0$ at point $A_0$. Point $A_1$ lies on circle $C_1$ so that $A_1$ is located $90^{\circ}$ counterclockwise from $A_0$ on $C_1$. Circle $C_2$ has radius $r^2$ and is internally tangent to $C_1$ at point $A_1$. In this way a sequence of circles $C_1,C_2,C_3,\ldots$ and a sequence of points on the circles $A_1,A_2,A_3,\ldots$ are constructed, where circle $C_n$ has radius $r^n$ and is internally tangent to circle $C_{n-1}$ at point $A_{n-1}$, and point $A_n$ lies on $C_n$ $90^{\circ}$ counterclockwise from point $A_{n-1}$, as shown in the figure below. There is one point $B$ inside all of these circles. When $r = \frac{11}{60}$, the distance from the center $C_0$ to $B$ is $\frac{m}{n}$, where $m$ and $n$ are relatively prime positive integers. Find $m+n$.


 
                    \begin{figure}[H]
                        \centering
                        \includegraphics[width=0.5\textwidth]{Problem/AIME_2017_II_12_Problem_Figure_1.png}
                    \end{figure}
                    


 
 \textbf{Answer: }110



\textbf{Solution 1}

Impose a coordinate system and let the center of $C_0$ be $(0,0)$ and $A_0$ be $(1,0)$. Therefore $A_1=(1-r,r)$, $A_2=(1-r-r^2,r-r^2)$, $A_3=(1-r-r^2+r^3,r-r^2-r^3)$, $A_4=(1-r-r^2+r^3+r^4,r-r^2-r^3+r^4)$, and so on, where the signs alternate in groups of $2$. The limit of all these points is point $B$. Using the sum of infinite geometric series formula on $B$ and reducing the expression: $(1-r)-(r^2-r^3)+(r^4-r^5)-...=\frac{1-r}{1+r^2}$, $(r-r^2)-(r^3-r^4)+(r^5-r^6)-...=\frac{r-r^2}{1+r^2}$. Thus, we get $B=\left(\frac{1-r}{r^2+1},\frac{r-r^2}{r^2+1}\right)$. The distance from $B$ to the origin is $\sqrt{\left(\frac{1-r}{r^2+1}\right)^2+\left(\frac{r-r^2}{r^2+1}\right)^2}=\frac{1-r}{\sqrt{r^2+1}}.$ Let $r=\frac{11}{60}$, and the distance from the origin is $\frac{49}{61}$.  $49+61=\boxed{110}$.

 





\textbf{Solution 2}

Let the center of circle $C_i$ be $O_i$. Note that $O_0BO_1$ is a right triangle, with right angle at $B$. Also, $O_1B=\frac{11}{60}O_0B$, or $O_0B = \frac{60}{61}O_0O_1$. It is clear that $O_0O_1=1-r=\frac{49}{60}$, so $O_0B=\frac{60}{61}\times\frac{49}{60}=\frac{49}{61}$. Our answer is $49+61=\boxed{110}$

 -william122

 





\textbf{Solution 3}

Note that there is an invariance, Consider the entire figure $\mathcal{F}$. Perform a $90^\circ$ counterclockwise rotation, then scale by $r$ with respect to $(1, 0)$. It is easy to see that the new figure $\mathcal{F}' \cup S^1 = \mathcal{F}$, so $B$ is invariant.

 Using the invariance, Let $B = (x,y)$. Then rotating and scaling, $B = (1-r(1+y), rx)$. Equating, we find $x = \frac{1-r}{r^2+1}, y = \frac{r-r^2}{r^2+1}$. The distance is thus $\frac{49}{61}$. Our answer is $49+61=\boxed{110}$

 -Isogonal

 





\textbf{Solution 4}

Using the invariance again as in Solution 3, assume $B$ is $d$ away from the origin. The locus of possible points is a circle with radius $d$. Consider the following diagram.
                    \begin{figure}[H]
                        \centering
                        \includegraphics[width=0.5\textwidth]{Solution/AIME_2017_II_12_Solution_Figure_1.png}
                    \end{figure}
                    

 Let the distance from $B$ to $(1,0)$ be $x$. As $B$ is invariant, $x = r(BB' + x) \implies x = r\frac{d\sqrt{2}}{1-r}$. Then by Power of a Point, $x(BB' + x) = (1-d)(1+d) \implies xr(BB' + x) = r(1-d)(1+d) \implies x^2 = r(1-d^2) \implies d^2 = \left(1 + \frac{2r}{(1-r)^2}\right)$. Solving, $d = \frac{49}{61}$. Our answer is $49+61=\boxed{110}$

 -Isogonal

 





\textbf{Solution 5 (complex)}

Let $A_0$ be the origin. Now note that the ratio of lengths of consecutive line segments is constant and equal to $r$. Now accounting for rotation by $\frac{\pi}{2}$ radians, we see that the common ratio is $ri$. Thus since our first term is $A_1=-r+ri$, the total sum (by geometric series formula) is $\frac{-r+ri}{1-ri}=\frac{-781+538i}{3721}$. We need the distance from $C_0=-1$ so our distance is $|B-C_0|=\sqrt{\left(\frac{-781}{3721}-(-1)\right)^2+\left(\frac{538i}{3721}\right)^2}=\sqrt{\frac{2401}{3721}}=\frac{49}{61}$. Our answer is $49+61=\boxed{110}$

 -chrisdiamond10

 By the way, \href{https://www.geogebra.org/calculator/wrtgmqmq}{here's} a Geogebra Diagram. ~r00tsOfUnity

 




	\item (\textit{AIME II 2017}) For each integer $n\geq3$, let $f(n)$ be the number of $3$-element subsets of the vertices of a regular $n$-gon that are the vertices of an isosceles triangle (including equilateral triangles). Find the sum of all values of $n$ such that $f(n+1)=f(n)+78$.


 
 \textbf{Answer: }245



\textbf{Solution 1}

Considering $n \pmod{6}$, we have the following formulas:

 $n\equiv 0$: $\frac{n(n-4)}{2} + \frac{n}{3}$

 $n\equiv 2, 4$: $\frac{n(n-2)}{2}$

 $n\equiv 3$: $\frac{n(n-3)}{2} + \frac{n}{3}$

 $n\equiv 1, 5$: $\frac{n(n-1)}{2}$

 To derive these formulas, we note the following:Any isosceles triangle formed by the vertices of our regular $n$-sided polygon $P$ has its sides from the set of edges and diagonals of $P$. Notably, as two sides of an isosceles triangle must be equal, it is important to use the property that same-lengthed edges and diagonals come in groups of $n$, unless $n$ is even when one set of diagonals (those which bisect the polygon) comes in a group of $\frac{n}{2}$. Three properties hold true of $f(n)$:

 When $n$ is odd there are $\frac{n(n-1)}{2}$ satisfactory subsets (This can be chosen with $n$ choices for the not-base vertex, and $\frac{n-1}{2}$ for the pair of equal sides as we have $n-1$ edges to choose from, and we must divide by 2 for over-count).*

 
                \begin{itemize}
                    \item Another explanation: For any diagonal or side of the polygon chosen as the base of the isosceles triangle, there is exactly 1 isosceles triangle that can be formed. So, the total number of satisfactory subsets is $\dbinom{n}{2}=\dfrac{n(n-1)}{2}.$
                \end{itemize}


                When $n$ is even there are $\frac{n(n-2)}{2}$ satisfactory subsets (This can be chosen with $n$ choices for the not-base vertex, and $\frac{n-2}{2}$ for the pair of equal sides as we have $n-1$ edges to choose from, one of them which is not satisfactory (the bisecting edge), and we must divide by 2 for over-count).

 When $n$ is a multiple of three we additionally over-count equilateral triangles, of which there are $\frac{n}{3}$. As we count them three times, we are two times over, so we subtract $\frac{2n}{3}$.

 Considering the six possibilities $n \equiv 0,1,2,3,4,5 \pmod{6}$ and solving, we find that the only valid solutions are $n = 36, 52, 157$, from which the answer is $36 + 52 + 157 = \boxed{245}$.

 





\textbf{Solution 2 (elaborates on the possible cases)}

In the case that $n\equiv 0\pmod 3$, there are $\frac{n}{3}$ equilateral triangles. We will now count the number of non-equilateral isosceles triangles in this case.

 Select a vertex $P$ of a regular $n$-gon. We will count the number of isosceles triangles with their vertex at $P$. (In other words, we are counting the number of isosceles triangles $\triangle APB$ with $A, B, P$ among the vertices of the $n$-gon, and $AP=BP$.)

 If the side $AP$ spans $k$ sides of the $n$-gon (where $k<\frac{n}{2}$), the side $BP$ must span $k$ sides of the $n$-gon, and, thus, the side $AB$ must span $n-2k$ sides of the $n$-gon. As $\triangle ABP$ has three distinct vertices, the side $AB$ must span at least one side, so $n-2k \ge 1$. Combining this inequality with the fact that $1\le k<\frac{n}{2}$ and $k\not = \frac{n}{3}$ (as $\triangle ABP$ cannot be equilateral), we find that there are $\lceil\frac{n}{2}\rceil-2$ possible $k$. 

 As each of the $n$ vertices can be the vertex of a given triangle $\triangle ABP$, there are $\left(\lceil \frac{n}{2} \rceil -2 \right)\cdot n$ non-equilateral isosceles triangles.

 Adding in the $\frac{n}{3}$ equilateral triangles, we find that for $n\equiv 0\pmod 3$: $f(n) = \frac{n}{3}+\left(\lceil \frac{n}{2} \rceil -2\right)\cdot n$.

 On the other hand, if $n\equiv 1, 2\pmod 3$, there are no equilateral triangles, and we may follow the logic of the paragraph above to find that $f(n)=\left(\lceil \frac{n}{2} \rceil -1\right)\cdot n$.

 We may now rewrite the given equation, based on the remainder $n$ leaves when divided by 3.

 
Case 1: $n\equiv 0\pmod 3$The equation $f(n+1)=f(n)+78$ for this case is $\left(\lceil \frac{n+1}{2} \rceil -1\right)\cdot (n+1)=\frac{n}{3}+\left(\lceil \frac{n}{2} \rceil -2\right)\cdot n+78$.

 In this case, $n$ is of the form $6k$ or $6k+3$, for some integer $k$.

 Subcase 1: $n=6k$Plugging into the equation above yields $k=6\rightarrow n=36$.

 Subcase 2: $n=6k+3$Plugging into the equation above yields $7k=75$, which has no integer solutions.

 Case 2: $n\equiv 1\pmod 3$The equation $f(n+1)=f(n)+78$ for this case is $\left(\lceil \frac{n+1}{2} \rceil -1\right)\cdot (n+1)=\left(\lceil \frac{n}{2} \rceil -1\right)\cdot n+78$.

 In this case, $n$ is of the form $6k+1$ or $6k+4$, for some integer $k$.

 Subcase 1: $n=6k+1$In this case, the equation above yields $k=8\rightarrow n=52$.

 Subcase 2: $n=6k+4$In this case, the equation above yields $k=26\rightarrow n=157$.

 Case 3: $n\equiv 2\pmod 3$The equation $f(n+1)=f(n)+78$ for this case is $\frac{n+1}{3}+\left(\lceil \frac{n+1}{2} \rceil -2\right)\cdot (n+1)=\left(\lceil \frac{n}{2} \rceil -1\right)\cdot n+78$.

 In this case, $n$ is of the form $6k+2$ or $6k+5$, for some integer $k$.

 Subcase 1: $n=6k+2$The equation above reduces to $5k=77$, which has no integer solutions.

 Subcase 2: $n=6k+5$The equation above reduces to $k=-80$, which does not yield a positive integer solution for $n$.

 In summary, the possible $n$ are $36, 52, 157$, which add to $\boxed{245}$.

 





\textbf{Solution 3}

We first notice that when a polygon has $s$ sides where $s\not\equiv 0\pmod{3}$, there cannot exist any three vertices that form an equilateral triangle. Also, the parity of $s$ and $s+1$ also matters, since they influence how many isosceles triangles including equilateral triangles exist in the polygon. We can model an equation $2x+y=s$, where the lines that are congruent connect the vertices that are $x$ vertices apart and the other line has vertices that are $y$ vertices apart. If $s$ is even, there are $\frac{s-2}{2}$ solutions for $(x,y)$ which would create a unique type of isosceles triangle. We subtract two since $y$ cannot be zero. If $s$ is odd, there are $\frac{s-1}{2}$ solutions for $(x,y)$.Next, we do casework on the congruence of $s\pmod{3}$ and the parody of $s$ using the information above:

 Case 1:

 $s\not\equiv 0\pmod{3}$, $s+1\not\equiv 0\pmod{3}$

 $s$ is even, $s+1$ is odd

 There are $\frac{s-2}{2}$ unique types of isosceles triangles in the polygon with $s$ sides. Each isosceles triangle has a unique point which connects the two congruent sides. Therefore, for each type of isosceles triangle, there exists $s$ of those triangles since the unique point can be any of the $s$ vertices. There are $\frac{s}{2}$ types of isosceles triangles in the polygon with $s+1$ sides and $s+1$ unique points for each type of triangle. Therefore, $\frac{s}{2}\cdot{(s+1)}-\frac{s-2}{2}\cdot{(s)}=78$. Solving, we get $s=52$.

 Case 2:

 $s\not\equiv 0\pmod{3}$, $s+1\not\equiv 0\pmod{3}$

 $s$ is odd, $s+1$ is even

 Using similar logic, there are $\frac{s-1}{2}\cdot{(s)}$ possible isosceles triangles in the $s$ sided polygon. There are $\frac{s-1}{2}\cdot{(s+1)}$ possible isosceles triangles in the $s+1$ sided polygon. The difference should be $78$, so $\frac{s-1}{2}\cdot{(s+1)}-\frac{s-1}{2}\cdot{(s)}=78$. Solving gives $s=157$.

 For both of these cases, we don't have to worry about equilateral triangles since in both cases $s,s+1\nmid3$.

 Case 3:

 $s\not\equiv 0\pmod{3}$, $s+1\equiv 0\pmod{3}$

 $s$ is even, $s+1$ is odd

 There are $\frac{s-2}{2}\cdot{(s)}$ possible isosceles triangles in the $s$ sided polygon. The $s+1$ case is a bit more complicated, as we have to consider equilateral triangles as well. In this case, there is one solution where $x=y$, which would produce an equilateral triangle. Therefore, we subtract that case to calculate only isosceles triangles with 2 congruent sides. Only isosceles triangles with exactly 2 congruent sides have a unique point, while there exist only $\frac{s+1}{3}$ distinct equilateral triangles in a polygon with $s+1$ sides, the rest are equivalent and symmetrical. Therefore, there are $(\frac{s}{2}-1)\cdot{(s+1)}+\frac{s+1}{3}$ isosceles triangles with at least 2 congruent sides in a polygon with $s+1$ sides. Putting it together, $\frac{s}{2}-1\cdot{(s+1)}+(\frac{s+1}{3})-\frac{s-2}{2}\cdot{(s)}=78$. Solving yields $5s=772$ which is impossible since $s$ has to be an integer. Therefore, this case is not valid.

 Case 4:

 $s\not\equiv 0\pmod{3}$, $s+1\equiv 0\pmod{3}$

 $s$ is odd, $s+1$ is even

 There are  $\frac{s-1}{2}\cdot{(s)}$ isosceles triangles in the $s$ sided polygon. Using the idea above, there are $(\frac{s-1}{2}-1)\cdot{(s+1)}$ isosceles triangles with 2 congruent sides in an $s+1$ sided polygon. There are $\frac{s+1}{3}$ equilateral triangles. Therefore,  $(\frac{s-1}{2}-1)\cdot{(s+1)}+(\frac{s+1}{3})-\frac{s-1}{2}\cdot{(s)}=78$. Looking at the left hand side, it is clear that $s$ has to be negative, which is not valid. Therefore, this case is not valid.

 Case 5:

 $s\equiv 0\pmod{3}$, $s+1\not\equiv 0\pmod{3}$

 $s$ is even, $s+1$ is odd

 There are a total of $(\frac{s-2}{2}-1)\cdot{(s)}+(\frac{s}{3})$ isosceles triangles in a polygon with $s$ sides. There are a total of $\frac{s}{2}\cdot{(s+1)}$ isosceles triangles in a polygon with $s+1$ sides. Therefore, $\frac{s}{2}\cdot{(s+1)}-(\frac{s-2}{2}-1)\cdot{(s)}-(\frac{s}{3})=78$. Simplifying, we get $s=36$.

 Case 6:

 $s\equiv 0\pmod{3}$, $s+1\not\equiv 0\pmod{3}$

 $s$ is odd, $s+1$ is even

 There are a total of $(\frac{s-1}{2}-1)\cdot{(s)}+(\frac{s}{3})$ isosceles triangles in the polygon with $s$ sides. There are $\frac{s-1}{2}\cdot{(s+1)}$ isosceles triangles in the polygon with $s+1$ sides. Therefore,$\frac{s-1}{2}\cdot{(s+1)}-(\frac{s-1}{2}-1)\cdot{(s)}-(\frac{s}{3})=78$. Simplifying, we get $7s=471$, which means $s$ is not an integer. Thus, this case is invalid.

 Adding all the valid cases, we obtain $52+157+36=\boxed{245}$

 





\textbf{Solution 4 (Bash)}

Notice that when $m$ is not a multiple of $3$, $f(m)=m\cdot \left\lfloor \frac{m-1}{2} \right\rfloor$, and when $m$ is a multiple of $3$, $f(m)=m\cdot \left\lfloor \frac{m-1}{2} \right\rfloor-\frac{m}{3}\cdot 2$.

 This is because we can choose any point on the $m$-gon to be the apex of the isosceles triangle, and we can choose from $\left\lfloor \frac{m-1}{2} \right\rfloor$ pairs of vertices that oppose each other with respect to the apex point. If there is an extra point that is exactly opposite to the first point we chose, $\frac{m-1}{2}$ is not an integer, and the extra point cannot form an isosceles triangle with the first point we chose (it would be a line). We take the floor to get rid of this case. 

 When $m$ is a multiple of $3$, the $m$-gon has $\frac{m}{3}$ equilateral triangles, but we have counted each thrice (once for each vertex), so we have to subtract 2 for each equilateral triangle, thus the $-\frac{m}{3}\cdot 2$ term.

 Now, define $g(n)=f(n+1)-f(n)$.

 We bash, hoping to find a pattern:

 \begin{alignat*}{3}f(3)  &{}= 3\cdot1 - 2\cdot1 = 3 - 2 &{}= 1  &\qquad g(3)  &{}= 4 - 1   &{}= 3 \\f(4)  &{}= 4\cdot1                  &{}= 4  &\qquad g(4)  &{}= 10 - 4  &{}= 6 \\f(5)  &{}= 5\cdot2                  &{}= 10 &\qquad g(5)  &{}= 8 - 10  &{}= -2 \\f(6)  &{}= 6\cdot2 - 2\cdot2 = 12 - 4 &{}= 8  &\qquad g(6)  &{}= 21 - 8  &{}= 13 \\f(7)  &{}= 7\cdot3                  &{}= 21 &\qquad g(7)  &{}= 24 - 21 &{}= 3 \\f(8)  &{}= 8\cdot3                  &{}= 24 &\qquad g(8)  &{}= 30 - 24 &{}= 6 \\f(9)  &{}= 9\cdot4 - 2\cdot3 = 36 - 6 &{}= 30 &\qquad g(9)  &{}= 40 - 30 &{}= 10 \\f(10) &{}= 10\cdot4                 &{}= 40 &\qquad g(10) &{}= 55 - 40 &{}= 15 \\f(11) &{}= 11\cdot5                 &{}= 55 &\qquad g(11) &{}= 52 - 55 &{}= -3 \\f(12) &{}= 12\cdot5 - 2\cdot4 = 60 - 8 &{}= 52 &\qquad g(12) &{}= 78 - 52 &{}= 26 \\f(13) &{}= 13\cdot6                 &{}= 78 &\qquad g(13) &{}= 84 - 78 &{}= 6 \\f(14) &{}= 14\cdot6                 &{}= 84 &\qquad g(14) &{}= 95 - 84 &{}= 11 \\f(15) &{}= 15\cdot7 - 2\cdot5 = 105 - 10 &{}= 95 &\qquad g(15) &{}= 112 - 95 &{}= 17 \\f(16) &{}= 16\cdot7                 &{}= 112 &\qquad g(16) &{}= 136 - 112 &{}= 24 \\f(17) &{}= 17\cdot8                 &{}= 136 &\qquad g(17) &{}= 132 - 136 &{}= -4 \\f(18) &{}= 18\cdot8 - 2\cdot6 = 144 - 12 &{}= 132 &\qquad g(18) &{}= 171 - 132 &{}= 39 \\f(19) &{}= 19\cdot9                 &{}= 171 &\qquad g(19) &{}= 180 - 171 &{}= 9 \\f(20) &{}= 20\cdot9                 &{}= 180 &\qquad g(20) &{}= 196 - 180 &{}= 16 \\f(21) &{}= 21\cdot10 - 2\cdot7 = 210 - 14 &{}= 196 &\qquad g(21) &{}= 220 - 196 &{}= 24 \\f(22) &{}= 22\cdot10                &{}= 220 &\qquad g(22) &{}= 253 - 220 &{}= 33 \\f(23) &{}= 23\cdot11                &{}= 253 &\qquad g(23) &{}= 248 - 253 &{}= -5 \\f(24) &{}= 24\cdot11 - 2\cdot8 = 264 - 16 &{}= 248 &\qquad g(24) &{}= 300 - 248 &{}= 52 \\f(25) &{}= 25\cdot12                &{}= 300 &\qquad g(25) &{}= 312 - 300 &{}= 12 \\f(26) &{}= 26\cdot12                &{}= 312 &\qquad g(26) &{}= 333 - 312 &{}= 21 \\f(27) &{}= 27\cdot13 - 2\cdot9 = 351 - 18 &{}= 333 &\qquad g(27) &{}= 364 - 333 &{}= 31 \\f(28) &{}= 28\cdot13                &{}= 364 &\qquad g(28) &{}= 406 - 364 &{}= 42 \\f(29) &{}= 29\cdot14                &{}= 406 &\qquad g(29) &{}= 400 - 406 &{}= -6 \\f(30) &{}= 30\cdot14 - 2\cdot10 = 420- 20 &{}= 400 \end{alignat*}

 We can see that every 6 rows there is a pattern. We will assume this pattern does not change for later numbers:

 \begin{aligned}g(3) &= 3, & g(4) &= 6, & g(5) &= -2, & g(6) &= 13, & g(7) &= 3, & g(8) &= 6, \\g(9) &= 10, & g(10) &= 15, & g(11) &= -3, & g(12) &= 26, & g(13) &= 6, & g(14) &= 11, \\g(15) &= 17, & g(16) &= 24, & g(17) &= -4, & g(18) &= 39, & g(19) &= 9, & g(20) &= 16, \\g(21) &= 24, & g(22) &= 33, & g(23) &= -5, & g(24) &= 52, & g(25) &= 12, & g(26) &= 21, \\g(27) &= 31, & g(28) &= 42, & g(29) &= -6\end{aligned}

 We can describe this pattern as \begin{aligned}g(6x+3) &= 7x+3 \\g(6x+4) &= 9x+6 \\g(6x+5) &= -x-2 \\g(6x+6) &= 13x+13 \\g(6x+7) &= 3x+3 \\g(6x+8) &= 5x+6 \\\end{aligned}

 where $x$ is a nonnegative integer.

 Equating these to $78$ yields\begin{aligned}g(6x+3) &= 7x+3 = 78 \implies x \not\in \mathbb{Z}\\g(6x+4) &= 9x+6 = 78 \implies x=8 \\g(6x+5) &= -x-2 = 78 \implies x \not\geq 0\\g(6x+6) &= 13x+13 = 78 \implies x=5\\g(6x+7) &= 3x+3 = 78 \implies x=25\\g(6x+8) &= 5x+6 = 78 \implies x \not\in \mathbb{Z}\\\end{aligned}

 Solving for the 3 cases where $x$ is a nonnegative integer yields $6\cdot8+4=52$, $6\cdot5+6=36$, and $6\cdot25+7=157$. Thus, our answer is $52+36+157=\boxed{245}$.

 


 




	\item (\textit{AIME II 2017}) A $10\times10\times10$ grid of points consists of all points in space of the form $(i,j,k)$, where $i$, $j$, and $k$ are integers between $1$ and $10$, inclusive. Find the number of different lines that contain exactly $8$ of these points.


 
 \textbf{Answer: }168



\textbf{Solution 1}

A line through the point $(a,b,c)$ must contain $(a \pm 1, b \pm 1, c \pm 1)$ on it as well, as otherwise, the line would not pass through more than 5 points. This corresponds to the 4 diagonals of the cube. With this information, we proceed with casework:

 $\textbf{Case 1}:$ The lines are not parallel to the faces

 We look at the one from $(1,1,1)$ to $(10,10,10)$. The lower endpoint of the desired lines must contain both a 1 and a 3, so it can be $(1,1,3), (1,2,3), (1,3,3)$. If $\textrm{min} > 0$ then the point $(a-1,b-1,c-1)$ will also be on the line for example, 3 applies to the other end.

 Accounting for permutations, there are $12$ ways, so there are $12 \cdot 4 = 48$ different lines for this case. 

 A nice visualization for this case is the following:

 
                    \begin{figure}[H]
                        \centering
                        \includegraphics[width=0.5\textwidth]{Solution/AIME_2017_II_14_Solution_Figure_1.png}
                    \end{figure}
                    

 The black vertex is the corner $(1,1,1)$. If we draw a line straight down through this point (black) you will hit $10$ points, if we draw a line straight down through the blue points you will hit $9$ points, and if we draw a line straight down through the red points you will hit $8$ points. Thus, for each pair of opposite vertices, we have $12$ lines corresponding to the $12$ red dots of length $8$, for a total of $12\cdot4=48$ points. 

 
$\textbf{Case 2}:$ The lines where the $x$, $y$, or $z$ is the same for all the points on the line. 

 WLOG, let the $x$ value stay the same throughout. Let the line be parallel to the diagonal from $(1,1,1)$ to $(1,10,10)$. For the line to have 8 points, the $y$ and $z$ must be 1 and 3 in either order, and the $x$ value can be any value from 1 to 10. In addition, this line can be parallel to 6 face diagonals. So we get $2 \cdot 10 \cdot 6 = 120$ possible lines for this case. 

 The answer is, therefore, $120 + 48 = \boxed{168}$

 Solution by stephcurry added to the wiki by Thedoge edited by Rapurt9 and phoenixfire

 





\textbf{Solution 2}

Look at one pair of opposite faces of the cube. There are $4$ lines say $l_1, l_2, l_3, l_4$ with exactly $8$ collinear points on the top face. For each of these lines, draw a rectangular plane that consists of one of the $l_i$ for $1 \leq i \leq 4$ and perpendicular to the top face. 

 There are $16$ lines in total on this plane. $10$ of which are parallel to one of the edges of the rectangular plane and $6$ of which are diagonals. There are $3$ pairs of opposite faces. So $3 \cdot 4 \cdot 16=192$ lines.

 But we are overcounting the lines of the diagonals of those rectangular planes twice. There are $4$ rectangular planes perpendicular to one pair of opposite faces. Thus $4 \cdot 6=24$ lines are overcounted.

 So the answer is $192-24=\boxed{168}$.

 Solution by phoenixfire

 





\textbf{Solution 3}

Considered the cases $(1, 2, ..., 8), (2, 3, ...,9), (3, 4, ..., 10)$ and reverse. Also, consider the constant subsequences of length 8 $(1, 1, ..., 1), (2, 2, ..., 2), ..., (10, 10, ..., 10)$. Of all the triplets that work they cannot be extended to form another point on the line in the $10 \times 10 \times 10$ grid but  we need to divide by 2 because reversing all the subsequences gives the same line. Thus the answer is 
 \[\frac{16^3 - 14^3 - 14^3 + 12^3}{2} = \boxed{168}\]

 





\textbf{Solution 4}

The lines can be defined as starting from $(a, b, c)$ with "slope" (vector) $(d, e, f)$. We impose the condition that at least one of $a - d, b - e, \textrm{ or } c - f$ is outside the range of $[1, 10]$ in order to ensure that $(a, b, c)$ is the first valid point on this line. Then, the line ranges from $(a, b, c), (a + d, b + e, c + f), \ldots, (a + 7d, b + 7e, c + 7f)$, where $1 \le a + 7d, b + 7e, c + 7f \le 10$, in which case at least one of $a + 8d, b + 8e, \textrm{ or } c + 8f$ is outside of the range $[1, 10]$ to ensure the line does not contain more than 8 points. For $1 \le a + 7d, b + 7e, c + 7f \le 10$ to be satisfied, the pairs $(a, d), (b, e), \textrm{ and } (c, f)$ can only be $(1, 0), (2, 0), \ldots (10, 0), (1, 1), (2, 1), (3, 1), (10, -1), (9, -1), \textrm{ and } (8, -1)$. Notice that there are only two pairs such that $a + 8d, b + 8e, \textrm{ or } c + 8f \not \in [1,  10]$, namely $(3, 1)$ and $(8, -1)$. Thus, our line must contain at least one of these two pairs. In addition, only the two pairs $(1, 1)$ and $(10, -1)$ satisfy $a - d, b - e, \textrm{ or } c - f \not \in [1, 10]$. Thus, we must also include at least one of these two pairs as well. Let us call these 4 pairs "important" pairs. Finally, we can include any valid pair for our third pair.

 Case 1: We repeat one of the important pairs

 There are 2 ways to choose from $(3, 1)$ and $(8, -1)$, and 2 ways to choose from $(1, 1)$ and $(10, -1)$. Then, there are 2 ways to choose the repeated pair. Next, we can arrange these pairs $3!/2! = 3$ ways. So, we have $2 \cdot 2 \cdot 2 \cdot 3 = 24.$

 Case 2: We use 3 distinct important pairs

 There are ${4 \choose 3}$ ways to choose the pairs (note that by pigeonhole principle, this guarantees we get at least one of each required pair). Then, there are $3! = 6$ ways to arrange it. We obtain $4 \cdot 6 = 24$.

 Case 3: We use 3 unique pairs, one that is not important.

 Again, there are $2 \cdot 2 = 4$ ways to choose the important pairs. Then, there are 12 ways to choose the non-important pair. Again, there are $3! = 6$ ways to arrange it. So, we get $4 \cdot 12 \cdot 6 = 288$.

 Summing all of it up and dividing by two (since we over counted each line and its reversal), we get $\frac{24 + 24 + 288}{2} = \boxed{168}.$

 


 


 




	\item (\textit{AIME II 2017}) Tetrahedron $ABCD$ has $AD=BC=28$, $AC=BD=44$, and $AB=CD=52$. For any point $X$ in space, define $f(X)=AX+BX+CX+DX$. The least possible value of $f(X)$ can be expressed as $m\sqrt{n}$, where $m$ and $n$ are positive integers, and $n$ is not divisible by the square of any prime. Find $m+n$.


 
 \textbf{Answer: }682



\textbf{Solution 2}

Set $a=BC=28$, $b=CA=44$, $c=AB=52$. Let $O$ be the point which minimizes $f(X)$.

 $\boxed{\textrm{Claim 1: } O \textrm{ is the gravity center } \ \frac {1}{4}(\vec A + \vec B + \vec C + \vec D)}$

 $\textrm{Proof:}$ Let $M$ and $N$ denote the midpoints of $AB$ and $CD$. From $\triangle ABD \cong \triangle BAC$ and $\triangle CDA \cong \triangle DCB$, we have  $MC=MD$, $NA=NB$ an hence $MN$ is a perpendicular bisector of both segments $AB$ and $CD$. Then if $X$ is any point inside tetrahedron $ABCD$, its orthogonal projection onto line $MN$ will have smaller $f$-value; hence we conclude that $O$ must lie on $MN$. Similarly, $O$ must lie on the line joining the midpoints of $AC$ and $BD$. $\square$

 $\boxed{\textrm{Claim 2: The gravity center } O \textrm{ coincides with the circumcenter.} \phantom{\vec A}}$

 $\textrm{Proof:}$  Let $G_D$ be the centroid of triangle $ABC$; then $DO = \frac 34 DG_D$ (by vectors). If we define $G_A$, $G_B$, $G_C$ similarly, we get $AO = \frac 34 AG_A$ and so on. But from symmetry we have $AG_A = BG_B = CG_C = DG_D$, hence $AO = BO = CO = DO$. $\square$

 Now we use the fact that an isosceles tetrahedron has circumradius $R = \sqrt{\frac18(a^2+b^2+c^2)}$. 

 Here $R = \sqrt{678}$ so $f(O) = 4R = 4\sqrt{678}$. Therefore, the answer is $4 + 678 = \boxed{682}$.

 





\textbf{Solution 3}


                    \begin{figure}[H]
                        \centering
                        \includegraphics[width=0.5\textwidth]{Solution/AIME_2017_II_15_Solution_Figure_1.png}
                    \end{figure}
                    Isosceles tetrahedron $ABCD$ or \href{https://en.wikipedia.org/wiki/Disphenoid}{Disphenoid} can be inscribed in a parallelepiped $AB'CD'C'DA'B,$ whose facial diagonals are the pares of equal edges of the tetrahedron $(AC = B'D',$ where $B'D' = BD).$ This parallelepiped is right-angled, therefore it is circumscribed and has equal diagonals. The center O of the circumscribed sphere (coincide with the centroid) has equal distance from each vertex. Tetrachedrons $ABCD$ and $A'B'C'D'$ are congruent, so point of symmetry O is point of minimum $f(X). f(O)= 4R$, where $R$ is the circumradius of parallelepiped.
 \[8R^2 = 2 CC'^2 = 2CD'^2 + 2D'B^2 + 2BC'^2,\] 
 \[2 CC'^2 = (CD'^2 + BC'^2) + (BC'^2 + BD'^2) + (CD'^2 + BD'^2) = AC^2 + AB^2+BC^2,\] 
 \[R = OC =\sqrt{\frac {AB^2 + AC^2 + AD^2}{8}}, f(O)= 4R = 4\sqrt {678}.\]

 \textbf{vladimir.shelomovskii@gmail.com, vvsss}   (Reconstruction)

 





\textbf{Solution 4}

Note that we can position the tetrahedron in a rectangular prism with one of the vertices $(x,y,z)$, due to the fact that the opposite sides are equal. Here is our setup:\begin{align*}A&=(0,0,0)\\B&=(x,y,0)\\C&=(x,0,z)\\D&=(0,y,z)\\\end{align*}
We obtain the system\begin{align*}x^2+y^2&=2704\\y^2+z^2&=1936\\z^2+x^2&=784\\\end{align*}
We see that $x^2+y^2+z^2=2712$. Let $X=(a,b,c)$, we have
 \[f(a,b,c)=\lvert  X-A\rvert+\lvert  X-B\rvert+\lvert  X-C\rvert+\lvert  X-D\rvert\]Taking the parital derivative and setting them equal to zero, we have\begin{align*}\frac{\partial f}{\partial a}=\frac{a}{\lvert  X-A\rvert}+\frac{a-x}{\lvert  X-B\rvert}+\frac{a-x}{\lvert  X-C\rvert}+\frac{a}{\lvert X-D\rvert}=0\\\frac{\partial f}{\partial b}=\frac{b}{\lvert  X-A\rvert}+\frac{b-y}{\lvert  X-B\rvert}+\frac{b}{\lvert  X-C\rvert}+\frac{b-y}{\lvert X-D\rvert}=0\\\frac{\partial f}{\partial c}=\frac{c}{\lvert  X-A\rvert}+\frac{c}{\lvert  X-B\rvert}+\frac{c-z}{\lvert  X-C\rvert}+\frac{c-z}{\lvert X-D\rvert}=0\end{align*}
Inspired by this symmetry, one may test the solution when $\lvert  X-A\rvert=\lvert  X-B\rvert=\lvert  X-C\rvert=\lvert  X-D\rvert$, which follows that
 \[a=\frac{x}{2},b=\frac{y}{2},c=\frac{z}{2}\]Indeed, the solution implies $AX=BX=CX=DX$, satisfying our hypothesis. We also know that $f$ concaves up, hence the minimal occurs in this regard.

 Therefore,
 \[f_{min}(X)=4\sqrt{(\frac{x}{2})^2+(\frac{y}{2})^2+(\frac{z}{2})^2}=4\sqrt{678}\]The requested sum is $4+678=\boxed{682}$

 




\end{enumerate}